\documentclass[11pt]{article}
\usepackage[a4paper,margin=29mm]{geometry}
\usepackage{amsmath,amssymb,amsthm,mathtools}
\usepackage{booktabs}
\usepackage{enumitem}
\usepackage{graphicx}
\usepackage{microtype}
\usepackage[hidelinks]{hyperref}

\newtheorem{theorem}{Theorem}[section]
\newtheorem{lemma}[theorem]{Lemma}
\newtheorem{proposition}[theorem]{Proposition}
\newtheorem{corollary}[theorem]{Corollary}
\newtheorem{conjecture}[theorem]{Conjecture}
\theoremstyle{remark}
\newtheorem{remark}[theorem]{Remark}

\newcommand{\SE}{\operatorname{SE}(2)}
\newcommand{\sn}{\operatorname{sn}}
\newcommand{\cn}{\operatorname{cn}}
\newcommand{\dn}{\operatorname{dn}}
\newcommand{\am}{\operatorname{am}}
\newcommand{\epsi}{\mathcal{E}}
\newcommand{\Exp}{\operatorname{Exp}}
\newcommand{\rk}{\operatorname{rank}}

\title{The Conjugate Locus of the Rotating Geodesics\\
in the Sub-Riemannian Problem on $\SE$}
\author{Igor Moiseev\thanks{\texttt{moiseev.igor@gmail.com}}}
\date{September 2026}

\begin{document}
\setlength{\emergencystretch}{2em}
\maketitle

\begin{abstract}
In the standard left-invariant sub-Riemannian problem on the group of plane
motions the conjugate times of a rotating geodesic are known: there is
one on each of two brackets in every Maxwell cell
$[p_n^1,p_{n+1}^1)$, $n\ge1$, the first conjugate time being the one on
the left bracket of cell~$1$.  These are two distinct times except at
$k=k_0$, $\psi\equiv K\pmod{2K}$, where they coincide at $2nK$.  We describe the image
of these points under the exponential map, the conjugate locus.  On every
open sheet the map has corank one with a kernel tangent to the level
sets of time, and the fold quantity along that kernel equals, up to an
explicit nonzero factor, one polynomial $Q$ in the Riccati variables of
the first-time problem; a positive rational certificate for $Q$ makes the
sheet a fold surface.  Its Maxwell edge and, below the critical modulus
$k_0$, its half-period edge are cusp edges; the common point at $k_0$ has
corank two and is a hyperbolic umbilic $D_4^+$; above $k_0$ the edge on
which the two phase branches meet is a fold.  The caustic is parametrized
by the midpoint amplitude and the modulus in an explicit half-angle chart
in which each of its eight open patches is embedded, the patches meet only
along common source seams, the small-modulus limit is the identity point,
and the separatrix limit leaves every compact set.  The proofs use the first
variation of the endpoint map, addition theorems, and the Riccati
identities of the conjugate-times paper.  The same identity governs every
later sheet: each open sheet is a fold surface; for $k\ne k_0$ the four
seams of a cell are a cusp edge at the Maxwell time, folds at the zeros
of $\beta_1$ and $\alpha_1$, and a cusp edge at the half period.  At
$k=k_0$ the $\alpha_1$ and half-period seams meet in a hyperbolic umbilic;
for large period count the fold quantity has an
explicit leading term of order $1/n$ whose sign is that of
$-\sn p\cn p$.  Which sheets meet is restricted by an octant lemma and
characterized through the scalar level function of the inverse problem:
a target lies on the left sheet of cell $n$ or the right sheet of cell
$j$ exactly when that function has a critical point at level $n-1$ or
$j$, with the residual argument above or below the half period; the
incidence map of the first four cells is computed numerically, the
modulus along the chart curve is proved unimodal.  A conditional sign criterion
implies at most one critical point at levels $\ge1$ and at most two
solutions per positive integer level on each admissible interval.
That criterion is verified by interval arithmetic on the endpoint domain
outside five declared exclusion sets ($2.56\%$ of the ambient volume;
the fifth is a low-modulus wedge on which the interval engine does not
close); the global count and
the types of the two maxima remain conjectural.  Exact identities and
selected singular values, kernels and edge derivatives have symbolic
and high-precision checks; Section~\ref{sec:numerics} records their
finite-sample scope and replay limitations.
\end{abstract}
\noindent\textit{Keywords:} sub-Riemannian geometry, group of motions of the plane, conjugate locus, Lagrangian singularities, computer-assisted proof.\\
\textit{MSC 2020:} 53C17, 58K05, 49K15, 65G20.

\section{Introduction}

The standard sub-Riemannian problem on $\SE$ is
\begin{equation}\label{eq:control-system}
 \dot x=u_1\cos\theta,\qquad \dot y=u_1\sin\theta,\qquad \dot\theta=u_2,
 \qquad \int_0^T(u_1^2+u_2^2)\,dt\to\min,
\end{equation}
with the identity as the initial point.  Its normal extremals are
parametrized by Jacobi elliptic functions, the pendulum structure of the
vertical part giving five strata $C_1,\dots,C_5$
\cite{MoiseevSachkov2010Maxwell}; the conjugate times, the cut time and the
optimal synthesis are known \cite{Sachkov2010Conjugate,Sachkov2011CutLocus},
and on the rotating stratum $C_2$ the first conjugate time, bracketed by
Sachkov, was made exact in
\cite{SE2FirstConjugate}: for every modulus $0<k<1$ and every initial phase
it is the unique zero of Sachkov's reduced Jacobian $J_2$ in the bracket
$[p_0,p_1]$, and it is reached by inverting one of two elliptic branches.

This paper describes the first conjugate locus itself: the surface
$\Sigma_1\subset\SE$ swept by the first conjugate points of the rotating
geodesics, as the image of a smooth sheet in the source under the
exponential map.  Its singularities are the singularities of a Lagrangian
map $\mathbb R^3\to\mathbb R^3$, so the expected local types are folds
($A_1$), cusp edges ($A_2$), swallowtails ($A_3$) and isolated umbilic
points ($D_4^\pm$) \cite{Arnold1972NormalForms}.  We use the map-germ indices $A_1$ (fold),
$A_2$ (cusp), $A_3$ (swallowtail) and $D_4^\pm$ (umbilics)
\cite{Whitney1955,GolubitskyGuillemin1973,Arnold1972NormalForms}; the
first sheet contains no swallowtail, since Theorems~\ref{thm:fold}
and~\ref{thm:edges} leave only folds and cusps on it.  For the umbilics we
follow Arnold's sign: $D_4^+$ (generating cubic $y_1^2y_2+y_2^3$, the
hyperbolic umbilic) is the class whose cubic has one real linear factor,
i.e. negative discriminant, and $D_4^-$ ($y_1^2y_2-y_2^3$, the elliptic
umbilic) the class with three; this is the convention of
Proposition~\ref{prop:d4}.  The results are six.

\begin{enumerate}[nosep,leftmargin=*]
\item \textbf{Fold.}  On the open first sheet $d\Exp$ has corank one with a
      kernel tangent to the level sets of time
      (Proposition~\ref{prop:kernel}), and the fold quantity along it is
      one polynomial $Q$ times an explicit nonzero factor
      (Theorem~\ref{thm:fold-identity}).  A positive rational certificate
      for $Q$ makes the whole open sheet a fold surface, including the
      tangency phases of \cite[Lemma~3.3]{SE2FirstConjugate}, where the
      time root has odd order at least three (Theorem~\ref{thm:fold}).
\item \textbf{Edges.}  The Maxwell edge is a cusp edge for every $k$.  For
      $k<k_0$ the half-period edge $p=2K$ is a second cusp edge; at
      $k=k_0$ its endpoint has corank two and is a hyperbolic umbilic
      $D_4^+$; for $k>k_0$ the edge $\alpha_1=0$, on which the two phase
      branches meet, is a fold (Theorem~\ref{thm:edges},
      Proposition~\ref{prop:d4}).
\item \textbf{Atlas.}  The caustic is parametrized by the midpoint
      amplitude $B$ and the modulus $k$, with an explicit target chart
      $(U,V,W)$ in which the immersion Jacobian is again $Q$
      (Proposition~\ref{prop:atlas}).
\item \textbf{Embedding.}  Each of the eight open patches is embedded; the
      patches and their seams meet only along common source seams; as
      $k\downarrow0$ the caustic collapses to the identity and as
      $k\uparrow1$ it leaves every compact set
      (Theorems~\ref{thm:univalence} and~\ref{thm:closed}).
\item \textbf{Every sheet.}  The kernel, the fold identity and the edge
      computations never use the first cell.  So every later sheet is a
      fold surface where $Q\ne0$, which the certificates give on every
      bracket; its seams are classified; the half-period seam degenerates
      to a hyperbolic umbilic at $k_0$ in every cell; and for large period
      count $Q$ has an explicit leading term (Theorem~\ref{thm:every-sheet},
      Proposition~\ref{prop:leading}), which makes the late-fold statement
      of the conjugate-times paper a corollary with the sign identified.
\item \textbf{Incidences.}  An octant lemma fixes which sheets can meet,
      the scalar level function of \cite{SE2Counts} characterizes the
      sheets as its integer critical levels, the incidence map of cells
      $1$--$4$ is computed numerically, the modulus along the chart curve is proved
      unimodal, the derivative of the level function is given in closed
      form, and a conditional sign criterion implies the count bound at
      levels $\ge1$.  After correction of two enclosure defects, its
      interval certificate was regenerated on the full $24\times20\times80$
      grid: the criterion holds on every cell outside five declared
      exclusion sets, the fifth being the low-modulus wedge
      $m\le0.06$, $\mu\ge0.94$, $0.41\le\Delta\varphi\le0.80$ together
      with one further cell, where the engine's budget passes leave
      boxes open (Section~\ref{sec:incidence}).
\end{enumerate}
Together with the exact conjugate times this is a complete description
of the first conjugate locus of $\SE$ in the rotating regime, its shape,
its singularities and its self-incidences, of the local structure of
every later sheet, and a partial account of incidences between sheets.
The count criterion is conditional and the global incidence and count
claims remain open.  Figure~\ref{fig:atlas} shows the atlas;
Figure~\ref{fig:strip} shows where the fold quantity degenerates.

The method is elementary where the original derivations used charts and
minors: the first variation of the endpoint map makes the
velocity transverse to the endpoint variations, so corank one reduces to
the nonvanishing of one derivative of the heading, and the kernel is the
null vector of the heading row.  Differentiating $J_2$ along that kernel
in raw coordinates gives the identity of Theorem~\ref{thm:fold-identity},
in which the transcendental Jacobi epsilon of the midpoint argument
cancels.  The edge classification is then a few lines each.  What is not
elementary is the positivity of $Q$: no short proof is known, and the
certificate of Section~\ref{sec:certificate} stands in for it.

\section{Coordinates and the first sheet}\label{sec:coordinates}

We use the notation of \cite{SE2FirstConjugate}.  Let $0<k<1$ be the
modulus, $m=k^2$, $k'=\sqrt{1-m}$, $K=K(k)$ and $E=E(k)$ the complete
integrals, and $\epsi(p)=\int_0^p\dn^2(s,k)\,ds$ Jacobi's epsilon.  In the
rotating coordinates
\begin{equation}\label{eq:coordinates}
 p=\frac t{2k},\qquad \tau=\psi+p,
\end{equation}
$\psi$ is the initial phase and $\tau$ the midpoint phase;
the terminal phase is $\psi_1=\psi+2p$. The source is
the cylinder $(\psi,k,t)$, $\psi\in\mathbb R/4K\mathbb Z$, for each of the two
rotation signs $s_2=\pm1$, and the energy is $H=1/2$.  The exponential map
is, for $s_2=+1$ \cite[Section 4.3]{MoiseevSachkov2010Maxwell},
\begin{equation}\label{eq:exp}
 \begin{aligned}
 x&=k\bigl[\dn\psi\,(\cn\psi-\cn\psi_1)+\sn\psi\,\Delta\bigr],\\
 y&=m\sn\psi\,(\cn\psi-\cn\psi_1)-\dn\psi\,\Delta,\\
 \theta&=\arcsin(k\sn\psi)-\arcsin(k\sn\psi_1),\\
 \psi_1&=\psi+2p,\qquad \Delta=2p+\epsi(\psi)-\epsi(\psi_1),
 \end{aligned}
\end{equation}
and $s_2=-1$ changes $(x,y,\theta)$ into $(-x,-y,\theta)$.  The heading
formula uses $(\arcsin(k\sn u))'=k\cn u$ and $\dot\theta=-\cn(\psi+t/k)$.

Sachkov's reduced Jacobian is
\begin{equation}\label{eq:j2}
 J_2=\alpha\sn^2\tau+\beta\cn^2\tau,\qquad
 \alpha=(1-m)\sn p\,\alpha_1,\qquad \beta=f_1\beta_1,
\end{equation}
\begin{equation}\label{eq:coeffs}
 \begin{aligned}
 \alpha_1&=\cn p\,\dn p\,[p-2\epsi(p)]+\sn p\{\dn^2p+\epsi(p)[p-\epsi(p)]\},\\
 f_1&=\cn p\,[\epsi(p)-p]-\dn p\,\sn p,\qquad
 \beta_1=\cn p\,\epsi(p)-\dn p\,\sn p,
 \end{aligned}
\end{equation}
and $\det d\Exp$ in the source coordinates $(t,\varphi,k)$, $\varphi=k\psi$, equals $J_2$
times the nonvanishing factor $-4k/[(1-m)(1-m\sn^2p\sn^2\tau)]$.  Let
$p_0=p_1^1(k)\in(K,2K)$ be the first positive zero of $f_1$ and
$p_1=\min\{2K,p_\alpha\}$, $p_\alpha$ the first zero of $\alpha_1$ after
$p_0$; $p_1=2K$ for $k\le k_0$ and $p_1=p_\alpha<2K$ for $k>k_0$, where
$k_0=0.9089085575\ldots$ is the modulus with $K=2E$.  On $(p_0,p_1)$ one has
$\alpha>0>\beta$, and with
\begin{equation}\label{eq:r}
 r=\frac{-\beta}{\alpha-\beta}\in(0,1),\qquad
 J_2=0\iff\sn^2\tau=r,
\end{equation}
the first-time theorem of \cite{SE2FirstConjugate} says that $r$ increases
strictly from $0$ to $1$ on $(p_0,p_1)$, and that the two phase branches
$\psi_\pm=\pm F(\arcsin\sqrt r,k)-p$ (mod $2K$) cover the phase circle
$\mathbb R/2K\mathbb Z$ exactly once.  Throughout we abbreviate
\begin{equation}\label{eq:abbrev}
 \begin{gathered}
 (s,c,d,e,q)=(\sn p,\cn p,\dn p,\epsi(p),p-\epsi(p)),\\
 (S,C,D)=(\sn\tau,\cn\tau,\dn\tau),\qquad \Delta_0=1-ms^2S^2.
 \end{gathered}
\end{equation}

\paragraph{The first sheet.}  The \emph{open first sheet} is
\[
 \Sigma=\{(\psi,k,t):\ p_0<p<p_1,\ \sn^2\tau=r(p,k)\},
\]
a smooth surface (Lemma~\ref{lem:sheet-smooth}).  Since $J_2$ depends on
$\psi$ modulo $2K$, it has two components over each $(p,k)$ modulo $2K$, one
per phase branch, and four on the source circle $\mathbb R/4K\mathbb Z$ (the
sign patches of Section~\ref{sec:atlas}).  We use the representatives
$\tau=u$ and $\tau=2K-u$, where $u=F(\arcsin\sqrt r,k)\in(0,K)$:
there $S=\sqrt r>0$ and $C$ has opposite signs on the two branches.
The other two patches are their shifts by $2K$, which negate both $S$
and $C$.  Its closure adds the \emph{Maxwell edge}
$p=p_0$ ($r=0$, $S=0$) and the \emph{upper edge} $p=p_1$ ($r=1$, $C=0$),
which for $k<k_0$ is $p=2K$ and for $k>k_0$ is $p=p_\alpha$; the two meet
at $k_0$.  The first caustic is $\Sigma_1=\Exp(\overline\Sigma)$.

\paragraph{Riccati variables.}  Following \cite[Lemma 3.2]{SE2FirstConjugate}
put
\begin{equation}\label{eq:ab}
 a=\frac{f_1}q,\qquad b=-\frac{\beta_1}e,\qquad z=ab=-\frac\beta\ell,\qquad
 \ell=eq,
\end{equation}
so that $a,b>0$ on $(p_0,p_1)$, $z=ab$ increases from $0$ to $1$, and
\begin{equation}\label{eq:riccati}
 a'=\frac{d^2+ca+ms^2a^2}{ds},\qquad
 z'=\frac{d^2b(1-z)+s^2a(1-m+mz)}{ds}>0 .
\end{equation}
The collapse identities $\beta=X-\ell$, $\alpha=(1-m)X$, $X=s\alpha_1$, give
\begin{equation}\label{eq:r-in-z}
 r=\frac z{(1-m)+mz},\qquad 1-r=\frac{(1-m)(1-z)}{(1-m)+mz},\qquad
 1-mr=\frac{1-m}{(1-m)+mz}.
\end{equation}
One more inequality will be used repeatedly: since $f_1=-cq-ds$,
\begin{equation}\label{eq:a-below-g}
 \begin{gathered}
 a=-c-\frac{ds}q<-c=:g,\\
 \text{hence}\qquad d^2+ca+ms^2a^2>d^2-c^2+ms^2a^2=s^2(1-m+ma^2)>0,
 \end{gathered}
\end{equation}
so $a'>0$ on $(p_0,p_1)$.

\begin{lemma}\label{lem:sheet-smooth}
$dJ_2\ne0$ on $\overline\Sigma$ minus the point $(p,\tau,k)=(2K,K,k_0)$.
\end{lemma}
\begin{proof}
$\partial_\tau J_2=2SCD(\alpha-\beta)$ is nonzero on $\Sigma$ ($S,C\ne0$,
$\alpha-\beta>0$).  At $S=0$, $\partial_pJ_2=\beta'C^2=f_1'\beta_1\ne0$
since $p_0$ is a simple zero of $f_1$ and $\beta_1(p_0)<0$.  At $C=0$,
$\partial_pJ_2=\alpha'$, which at $p_\alpha$ is $(1-m)s\alpha_1'\ne0$ and
at $p=2K$ is $(1-m)(-1)\alpha_1(2K)=2(1-m)(K-2E)$, zero exactly at $k_0$.
\end{proof}

\section{Corank one and the kernel}\label{sec:kernel}

Write $\eta_\psi,\eta_k\in T_q\SE$ for the endpoint variations at fixed $t$
and $\xi=\dot g$ for the velocity, $q=\Exp(\psi,k,t)$, and let $\lambda(t)$
be the covector of the extremal at $q$.

\begin{lemma}[First variation]\label{lem:first-variation}
$\langle\lambda(t),\eta_\psi\rangle=\langle\lambda(t),\eta_k\rangle=0$ and
$\langle\lambda(t),\xi\rangle=2H=1$.  Consequently
\[
 \rk d\Exp=\rk[\eta_\psi,\eta_k]+1 .
\]
\end{lemma}
\begin{proof}
Let $\phi_t$ be the Hamiltonian flow and $\varsigma=\lambda\,dq$ the Liouville
form, with the symplectic form $\omega=d\varsigma=d\lambda\wedge dq$ and
$X_H$ defined by $\iota_{X_H}\omega=-dH$, so that $\dot q=\partial_\lambda H$,
$\dot\lambda=-\partial_qH$.  Since $H$ is quadratic in the fiber, $\langle\lambda,\partial_\lambda H\rangle=2H$
and $\mathcal L_{X_H}\varsigma=d(2H)-dH=dH$, so $\phi_t^*\varsigma=\varsigma+t\,dH$.
Restricted to the fiber $T^*_e\SE$, where $\varsigma=0$, this reads
$\langle\lambda(t),d\Exp_t(\delta\lambda)\rangle=t\,dH(\delta\lambda)$ for
every fiber variation $\delta\lambda$
(see also \cite[Ch.~8]{AgrachevBarilariBoscain2020}).  Our family lies on
$H=1/2$, so the right side vanishes for $\delta\lambda=\partial_\psi,\partial_k$;
and $\langle\lambda,\xi\rangle=\langle\lambda,\partial_\lambda H\rangle=2H$.
Thus $\xi\notin\ker\lambda(t)\supset\operatorname{span}(\eta_\psi,\eta_k)$.
\end{proof}

\begin{proposition}[Kernel]\label{prop:kernel}
At a conjugate point $d\Exp$ has corank one if and only if
$(\eta_\psi,\eta_k)\ne(0,0)$, and then its kernel is spanned by a vector
with zero time component, namely by $v=(F_k,-F_\psi,0)$ for any target
coordinate $F$ with $(F_\psi,F_k)\ne(0,0)$.  On the open first sheet
\begin{equation}\label{eq:theta-psi}
 \theta_\psi=k\bigl[\cn\psi-\cn(\psi+2p)\bigr]
 =\frac{2k\,sSdD}{\Delta_0}\ne0,
\end{equation}
so $d\Exp$ has corank one on all of $\Sigma$, including the tangency
phases and the modulus $k_0$, with kernel
$v=(\theta_k,-\theta_\psi,0)$.
\end{proposition}
\begin{proof}
If $v_\psi\eta_\psi+v_k\eta_k+v_t\xi=0$, pairing with $\lambda(t)$ gives
$v_t=0$ by Lemma~\ref{lem:first-variation}; so the kernel consists of the
$(v_\psi,v_k,0)$ with $v_\psi\eta_\psi+v_k\eta_k=0$, it is nonzero at a
conjugate point because $\det[\eta_\psi,\eta_k,\xi]=0$ forces
$\eta_\psi\parallel\eta_k$, and it is one-dimensional exactly when the pair
is not zero.  Any row
$F$ with nonzero gradient then gives the null vector.  The heading
derivative is $\theta_\psi=k[\cn\psi-\cn(\psi+2p)]$ by \eqref{eq:exp}, and
the addition theorem $\cn(\tau\mp p)=(Cc\pm SsDd)/\Delta_0$ gives
\eqref{eq:theta-psi}; on $\Sigma$, $s=\sn p>0$ for $K<p<2K$ and
$|S|=\sqrt r>0$.
\end{proof}

\begin{remark}
The kernel is tangent to the level sets of physical time.  This is the
general fact behind Lemma~\ref{lem:first-variation}: at a conjugate time
the kernel of $d\Exp$ is the kernel of $d\Exp_t$ on the energy level, since
$\Exp$ is an immersion in $t$.  It explains why the kernel in the
half-angle chart is the cross product $dt\times dW$ of the time and the
half-angle differentials.
\end{remark}

\section{The fold identity}\label{sec:fold}

Let $Q$ be the polynomial of nineteen terms
\begin{equation}\label{eq:Q}
 \begin{aligned}
 Q(a,b,c,s,m)={}&-3a^3b^2cm-a^2b^3cm+3a^2b^2ms^2-8a^2b^2m+3a^2b^2\\
 &+2a^2bcm+a^2ms^2-a^2s^2+ab^3ms^2-ab^3-2ab^2cm\\
 &-5abms^2+8abm+3abs^2-6ab+acm-ac+3bcm-3bc,
 \end{aligned}
\end{equation}
evaluated at the Riccati variables \eqref{eq:ab} and $(c,s,m)=(\cn p,\sn p,k^2)$.

\begin{theorem}[Fold identity; computer-assisted]\label{thm:fold-identity}
On the open first sheet, along the kernel $v=(\theta_k,-\theta_\psi,0)$ and
at fixed $t$, the following identity holds; its final reduction is a
computer-algebra verification in the Jacobi ring (see the proof).
\begin{equation}\label{eq:fold-identity}
 v(J_2)=-\,\frac{2(1-m)\,q^4e^4\,S}{s\,d\,D\,\Delta_0\,(\alpha-\beta)^2}\;Q .
\end{equation}
Apart from the displayed minus sign and the signed factor $S$, every
factor multiplying $Q$ is positive.  Hence a point of $\Sigma$ is a fold
of $\Exp$ if and only if $Q\ne0$, and $v(\det d\Exp)$ has the sign of
$-SQ$ times the sign of the prefactor in \eqref{eq:j2}.  On the chosen
$S>0$ representatives this reduces to the sign of $-Q$ times that
prefactor; the phase-shift symmetries preserve the fold type.
\end{theorem}
\begin{proof}
In the coordinates $(\tau,p,k)$ of \eqref{eq:coordinates},
$\partial_\psi|_{k,t}=\partial_\tau$ and
$\partial_k|_{\psi,t}=\partial_k|_{\tau,p}-(p/k)(\partial_\tau+\partial_p)$,
so
\[
 v(J_2)=\theta_k\,\partial_\tau J_2
 -\theta_\psi\bigl[\partial_kJ_2-(p/k)(\partial_\tau J_2+\partial_pJ_2)\bigr].
\]
Here $\partial_\tau J_2=2SCD(\alpha-\beta)$, $\partial_pJ_2$ follows from
$s'=cd$, $c'=-sd$, $d'=-msc$, $e'=d^2$, $q'=ms^2$, and $\partial_k$ at fixed
argument uses the classical modulus derivatives
\begin{gather*}
 \partial_k\sn=\frac{ms c^2-\gamma cd}{k(1-m)},\qquad
 \partial_k\cn=\frac{s(\gamma d-msc)}{k(1-m)},\\
 \partial_k\dn=\frac{ks(\gamma c-sd)}{1-m},\qquad
 \partial_k\epsi=\frac{k(scd-c^2e-(1-m)ps^2)}{1-m},
\end{gather*}
with $\gamma=\epsi(u)-(1-m)u$ at the argument $u$.  The heading derivative
at fixed $(\psi,t)$ is, from \eqref{eq:exp} and $(s+k\partial_ks)/d=(sd-\gamma c)/(1-m)$,
\begin{equation}\label{eq:theta-k}
 \theta_k=\frac{s_0d_0-s_1d_1+(\gamma_1-\gamma_0)c_1}{1-m}
 -\frac{\gamma_0\,\theta_\psi}{k(1-m)}+2pc_1,
\end{equation}
where $(s_0,c_0,d_0)$ and $(s_1,c_1,d_1)$ are the Jacobi functions at
$\tau-p$ and $\tau+p$ (in this proof only), given by the addition theorems, and $\gamma_0,\gamma_1$
the values of $\gamma$ there; the addition theorem for $\epsi$ gives
$\gamma_1-\gamma_0=2\gamma_p-2msS^2cd/\Delta_0$ and
$\gamma_0=\gamma_\tau-\gamma_p+mSs\,s_0$ with $\gamma_p=e-(1-m)p$.  The
midpoint quantity $\gamma_\tau=\epsi(\tau)-(1-m)\tau$ enters $\theta_k$
and $\partial_k\sn\tau$ with opposite signs and cancels identically in
$v(J_2)$, which is therefore a rational function of
$(s,c,d,e,q,S,C,D,m,k)$ with denominator $\Delta_0^2k(1-m)CD$ and a
numerator of $155$ terms in normal form modulo $c^2=1-s^2$, $d^2=1-ms^2$,
$C^2=1-S^2$, $D^2=1-mS^2$, $k^2=m$.  Substituting $S^2=r$ from \eqref{eq:r}
leaves the single monomial $SC$, so that
$v(J_2)=SP_\Sigma/[(\alpha-\beta)^3\Delta_0^2k(1-m)D]$ with $P_\Sigma$ a
polynomial in the $p$-variables, and with
$\Delta_0(\alpha-\beta)=\alpha-d^2\beta$ the claim \eqref{eq:fold-identity}
becomes the polynomial identity
\[
 P_\Sigma\,s\,d=-2kqe(1-m)^2(\alpha-d^2\beta)\,q^3e^3Q,
\]
which holds exactly in the ring
(\path{research/e74_step2_fold_identity.py}; the script also checks
$\theta_k$, $\theta_\psi$ and $v(J_2)$ against finite differences of
\eqref{eq:exp} to $3\times10^{-20}$).  Positivity of the factors:
$1-m,q,e,d,D,\Delta_0>0$, $s>0$ on $(K,2K)$,
$|S|=\sqrt r>0$, and $(\alpha-\beta)^2>0$.  On the two chosen
representatives $S>0$; on their $2K$ shifts the identity carries the
corresponding negative sign of $S$.  Finally $\det d\Exp=FJ_2$ with $F\ne0$ and
$v(FJ_2)=Fv(J_2)$ on $J_2=0$.
\end{proof}

\section{The certificate and the fold theorem}\label{sec:certificate}

The polynomial $Q$ is positive on the open first sheet.  No short proof is
known, and the positivity rests on an exact certificate; its variables are
set in sans-serif to keep them apart from the coordinates.

\begin{theorem}[Fold-only open first sheet]\label{thm:fold}
$Q>0$ on the open first sheet and on its edge $\alpha_1=0$ for $k>k_0$.  Consequently
the open first sheet, including the tangency phases, consists
of folds of $\Exp$, and so does the edge $\alpha_1=0$ for $k>k_0$.
\end{theorem}
\begin{proof}
With $g=-c$, $\rho=sd$, $\mathsf u=g-a=\rho/q$, $\mathsf v=b-g=\rho/e$,
$L_0=(1-m)/m$, the bracket inequalities give $0<\mathsf u<g<1$, $\mathsf v>0$,
$ab\le1$, and
\[
 me-(1-m)q=m\int_0^p\cn^2>0,\qquad e>E>1,
\]
so $L_0\mathsf v<\mathsf u$ and $\mathsf v<\rho$.  Put $\mathsf x=\mathsf u/g$,
$\mathsf t=L_0\mathsf v/\mathsf u$, $\mathsf y=g^2$; then
$0<\mathsf x,\mathsf t,\mathsf y<1$, and $\mathsf v<\rho$, $ab\le1$ read
\begin{equation}\label{eq:BZ}
 \mathsf B=L_0^2(1-\mathsf y)(L_0+\mathsf y)-(1+L_0)\mathsf x^2\mathsf t^2\mathsf y>0,\qquad
 \mathsf Z=L_0-\mathsf y(1-\mathsf x)(L_0+\mathsf x\mathsf t)\ge0 ,
\end{equation}
with $\mathsf Z=0$ exactly at $ab=1$, i.e. at $\alpha_1=0$ by \eqref{eq:ab}.
Direct substitution gives
\[
 Q=\frac{g^2\mathsf x}{L_0^3(1+L_0)}\,P(L_0,\mathsf x,\mathsf t,\mathsf y),
\]
and exact rational arithmetic verifies
\begin{equation}\label{eq:certificate}
 P=S_0+\mathsf B\,S_B+\mathsf Z\,S_Z,
\end{equation}
where $S_0,S_B,S_Z$ are sums of products $L_0^{i_0}\mathcal B^{n_1}_{i_1}(\mathsf x)
\mathcal B^{n_2}_{i_2}(\mathsf t)\mathcal B^{n_3}_{i_3}(\mathsf y)$ of Bernstein basis
functions \cite{Bernstein1915,Handelman1988,Farouki2012} with $216$
strictly positive rational coefficients, of degree
tuples at most $(5,5,4,3)$, $(2,3,2,1)$, $(4,3,3,2)$ in $(L_0,\mathsf x,\mathsf t,\mathsf y)$.  The sparse table
is \path{data/open_sheet_q_certificate.json} (its digest is given in the
code and data section); it has been re-verified independently with
exact fractions only (\path{scripts/review/verify_q_certificate_fractions.py}).
In the open cube every Bernstein function is positive, so $S_0>0$, and
$\mathsf B>0$, $\mathsf Z\ge0$ give $P>0$, including at $\mathsf Z=0$.  Theorem~\ref{thm:fold-identity}
concludes; on the edge $\alpha_1=0$ the point is an interior point of the
smooth singular surface (Lemma~\ref{lem:sheet-smooth}) where the two phase
branches meet ($C=0$, $S=1$), Proposition~\ref{prop:kernel} applies, and
\eqref{eq:fold-identity} holds with $S=1$: both sides are real-analytic
on the smooth surface $\{J_2=0\}$ through the edge
(Lemma~\ref{lem:sheet-smooth}), and the edge is nowhere dense in it, so
the identity proved on the open sheet persists there.
\end{proof}

\begin{figure}[t]
\centering
\includegraphics[width=\textwidth]{../figures/caustic_fold_strip.pdf}
\caption{The fold quantity on the source rectangle of the first sheet,
$z=\sin^2B=r$ against $k$: $\log_{10}Q$ (color and white contour labels).
$Q$ is positive inside and tends to zero at the edge $p=2K$ for
$k\le k_0$, where the kernel derivative of $J_2$ must vanish because that
edge is a cusp edge, and along the separatrix $k\to1$.  At
the Maxwell edge $z=0$ it is the factor $S=\sqrt z$ in
\eqref{eq:fold-identity}, not $Q$, that vanishes.  The $\alpha_1$ edge above
$k_0$ carries $Q>0$: it is a fold.}
\label{fig:strip}
\end{figure}

\section{The edges}\label{sec:edges}

For a corank-one map with determinant $\mathcal J$, smooth singular set
and a smooth kernel extension $v$ near a singular point, $v(\mathcal J)\ne0$
is the fold criterion, and $v(\mathcal J)=0$, $v^2(\mathcal J)\ne0$,
$d\mathcal J\wedge d(v\mathcal J)\ne0$ is the criterion for a cusp $A_2$
(Morin's criterion, \cite[Appendix~A, Theorem~A.1]{SajiUmeharaYamada2008Recognition}).
Here $A_1$ (fold) and $A_2$ (cusp) use the Morin map convention;
the corresponding generating-function types are $A_2$ and $A_3$.
We use $\mathcal J=J_2$, which differs from
$\det d\Exp$ by a nonvanishing factor, harmless for both criteria.

\begin{theorem}[Edges]\label{thm:edges}
\begin{enumerate}[nosep,leftmargin=*]
\item[(i)] On the Maxwell edge $\{S=0,\ f_1=0\}$, for every $0<k<1$,
      $\eta_\psi=0$, the kernel is the phase direction, and the point is an
      $A_2$ cusp.
\item[(ii)] At $(p,\tau)=(2K,K)$, $\eta_\psi=0$, $d\theta=0$, and
      \[
       \eta_k=2(e-q)\Bigl(1,-\frac k{k'},0\Bigr),\qquad e=2E,\ q=2(K-E).
      \]
      Hence $d\Exp$ has corank one for $k\ne k_0$ and corank two at $k_0$.
      For $k\ne k_0$ the point is an $A_2$ cusp; for $k<k_0$ it is the
      upper edge of the first sheet.
\item[(iii)] For $k>k_0$ the edge $\alpha_1=0$ is a fold
      (Theorem~\ref{thm:fold}).
\end{enumerate}
\end{theorem}
\begin{proof}
(i) At $S=0$ one has $\Delta_0=1$, $D=1$, $C=\pm1$, and from \eqref{eq:exp}
by the addition theorems
\[
 x_\psi=-2kC\,d\,f_1,\qquad y_\psi=2ms\,f_1,\qquad \theta_\psi=0,
 \qquad \theta_k=\frac{2C\beta_1}{1-m};
\]
so on the edge $\eta_\psi=0$, and $\theta_k\ne0$ because $\beta_1<0$ on
$(K,b_1)$: corank one, kernel $\partial_\psi$, by Proposition~\ref{prop:kernel}.
The extension $v=(\theta_k,-\theta_\psi,0)$ is smooth through the edge and
equals the kernel there.  On the sheet $v(J_2)=SG$ by
\eqref{eq:fold-identity}, with $G\ne0$ at the edge: $f_1=0$ gives $a=0$ and
$Q=-3(1-m)bc>0$.  So $v(J_2)=0$ at the edge, and $v$ is tangent to the
sheet there.  A derivative along a vector tangent to the sheet depends only
on the restriction to the sheet, hence
\[
 v^2(J_2)=dS(v)\,G=\frac{2\beta_1}{1-m}\,G
 =-\frac{12(1-m)q^4e^3\beta_1^2\,c}{s\,d\,\alpha^2}>0,
\]
using $dS(v)=\theta_kCD=2\beta_1/(1-m)$ and, at the edge, $\beta=0$,
$\Delta_0=D=1$, $Q=3(1-m)c\beta_1/e$.  The same tangent derivative shows
$d(vJ_2)(v)\ne0$ while $dJ_2(v)=0$, which is the rank condition.

(ii) At $(2K,K)$: $s=0$, $c=-1$, $d=1$, $C=0$, $S=1$, $D=k'$, $\theta=0$.
From \eqref{eq:exp}, $x=-2kqS$ and $y=-2qD$ along $s=0$, so $x_\psi=-2kqCD=0$
and $y_\psi=2mqSC=0$; $\theta_\psi=0$ by \eqref{eq:theta-psi} and
$\theta_k=0$ by \eqref{eq:theta-k}; $\dot\theta=-\cn(3K)=0$.  The modulus
derivatives at fixed $(\psi,t)$ are $x_k=2(e-q)$ and $y_k=2k(q-e)/k'$
(\path{research/e74_step3_edges.py}, exact).  Since $e-q=2(2E-K)$
vanishes exactly at $k_0$, Lemma~\ref{lem:first-variation} gives the
corank statement.  For $k\ne k_0$ take the extension $v=(x_k,-x_\psi,0)$
from the $x$ row.  At the point $v(J_2)=x_k\partial_\tau J_2=0$ because
$C=0$, and $v$ is tangent to the sheet ($\partial_\tau J_2=0$,
$\partial_pJ_2\ne0$).  Differentiating once more along $v=x_k\partial_\tau$
at the point,
\[
 v^2(J_2)=x_k\bigl[x_k\,\partial_\tau^2J_2-x_{\psi\psi}\,\partial_k^{\rm src}J_2\bigr]
 =-12eq(1-m)(e-q)^2\ne0,
\]
from $\partial^2_\tau J_2=-2eq(1-m)$, $x_{\psi\psi}=2kq(1-m)$ and
$\partial^{\rm src}_kJ_2=(e-q)e/k$, where $\partial^{\rm src}_k$ is the
$k$-derivative at fixed $(\psi,t)$.  The bracket, $-6eq(1-m)(e-q)$, is
$d(vJ_2)(\partial_\tau)$ while $dJ_2(\partial_\tau)=0$: the rank
condition.  For $k<k_0$ the point is the upper end $p=p_1=2K$ of the first
sheet; for $k>k_0$ it lies on the closure of the right sheet of cell~$1$,
$\Sigma_1^+$ in the notation of Section~\ref{sec:every-sheet}.
\end{proof}

\begin{proposition}[Hyperbolic umbilic; computer-assisted]\label{prop:d4}
At $k=k_0$ the point $(p,\tau)=(2K,K)$ is a $D_4^+$ Lagrangian singularity
of $\Exp$.
\end{proposition}
\begin{proof}
$\Exp$ is Lagrangian: homogeneity identifies $(\lambda,t)$ on the energy
cylinder with $t\lambda$ in the fiber, the flow maps the fiber to a
Lagrangian submanifold, and $\Exp$ is its projection.  For the germ we use
the amplitude chart $(A,B,k)=(\am p,\am\tau,k)$ and the half-angle target
chart $(U,V,W)$ of Proposition~\ref{prop:atlas} below.  Here $e=\epsi(2K)=K(k_0)>1$ and $k'=\sqrt{1-k_0^2}$.  The vectors
$(1/k,0,0)$, $(0,1,0)$, $(-e,0,k(1-k^2))$ have $dV=1$ on the first while the
last two span the kernel, and the image of $d\Exp$ is the $V$-axis, so
$(U,W)$ are cokernel coordinates.  Hessian contractions of the chart
formulas \eqref{eq:atlas-target}, computed by computer algebra and
confirmed numerically to $10^{-10}$, after straightening
the regular coordinate $\sigma=V-V_0$ and removing pure $\sigma^2$ target
terms, give the singular quadratic components in kernel coordinates
$(y_1,y_2)$
\[
 U_2=\mathsf a y_1^2+\mathsf b y_2^2+\mathsf c\,\sigma y_2,\qquad
 W_2=\mathsf d y_1y_2+\mathsf f\,\sigma y_1,
\]
\[
 \mathsf a=\frac{ek^2}{2k'},\quad \mathsf b=\frac{ek^2k'(e^2-1)}2,\quad
 \mathsf c=-e^2kk',\quad \mathsf d=ek,\quad \mathsf f=-1 .
\]
Since $\Exp$ is Lagrangian, the intrinsic second derivative
$\ker\to\operatorname{coker}$ is the gradient of a binary cubic under the
symplectic pairing of the kernel with the cokernel; let $(m_{ij})$ be that
pairing, which turns $(U_2,W_2)$ into a gradient in $(y_1,y_2)$ (it is
unique up to scale, and the sign of the discriminant is scale-invariant);
Hessian symmetry forces $m_{12}\mathsf d=2m_{21}\mathsf a$,
$m_{22}\mathsf d=2m_{11}\mathsf b$, $m_{11}\mathsf c=m_{22}\mathsf f$, and
$\mathsf c-2\mathsf b\mathsf f/\mathsf d=-kk'\ne0$ gives $m_{11}=m_{22}=0$.
With $m_{21}=1$ the generating-family jet is
\[
 \widehat F_3=\mathsf a y_1^2y_2+\frac{\mathsf b}3y_2^3
 +\sigma\Bigl(\frac{\mathsf a\mathsf f}{\mathsf d}y_1^2+\frac{\mathsf c}2y_2^2\Bigr),
\]
whose binary cubic at $\sigma=0$ has discriminant
$-4\mathsf a^3\mathsf b/3=-e^4k^8(e^2-1)/(12k'^2)<0$: the $D_4^+$ sign
\cite[Ch.~5, \S2.3, p.~97]{ArnoldGivental1990Symplectic}.
The Jacobi algebra is spanned by
$1,y_1,y_2,y_1^2$; the two cokernel coordinates represent $y_1,y_2$ and
the regular parameter represents $y_1^2$ with
coefficient $\mathsf a(\mathsf f/\mathsf d-\mathsf c/(2\mathsf b))=k/[2k'(e^2-1)]>0$,
so the three-parameter deformation is versal
\cite[Ch.~5, \S2.5, pp.~98--99]{ArnoldGivental1990Symplectic}.
For completeness, the cubic is $3$-right-determined.  After positive
linear rescaling it is $f=x^2y+y^3$.  The five polynomials
$x^2f_x,xyf_x,y^2f_x,x^2f_y,y^2f_y$ form a basis of the quartic
homogeneous polynomials.  If $R$ vanishes to order four, the same five
expressions for $f_t=f+tR$ have an invertible coefficient matrix over
the quartic monomials near the origin, uniformly for $0\le t\le1$.
Thus $R=df_t(X_t)$ for a smooth vector field $X_t$ vanishing to order
two.  The local flow of $-X_t$ fixes the origin and satisfies
$f_t\circ\Phi_t=f$, proving the assertion.  Together with versality,
this identifies the full germ, by the generating-family equivalence
theorem \cite[Ch.~5, \S1.4, p.~93]{ArnoldGivental1990Symplectic}.
The same three local types occur on the ellipsoid \cite{ItohKiyohara2021}.
\end{proof}

\begin{remark}[Intrinsic check]\label{rem:intrinsic}
The real form of the generating cubic can be checked intrinsically;
classification of the full germ still requires the versality established
above.  On the kernel plane
$\ker d\Exp=\operatorname{span}(\partial_\psi,\partial_k)$ the second
differential modulo the image $\mathbb R\xi$ is a quadratic map
$\Phi:\ker\to\operatorname{coker}$, which for a Lagrangian map is the
gradient of the generating cubic in suitable linear coordinates; the
quadratic form $v\mapsto\det D\Phi(v)$ is then the Hessian form of that
cubic up to a constant, definite for $D_4^-$ and indefinite for $D_4^+$.
At $40$ digits its discriminant at our point is $2.2\times10^5>0$
(\path{research/e74_step4_d4.py}), in agreement with the proposition.
This discriminant alone is insufficient: the Lagrangian maps
$(\sigma,2xy,x^2+3y^2)$ and $(\sigma,2xy+2\sigma x,x^2+3y^2)$
have the same intrinsic quadratic map at the origin, but their rank-one
loci are respectively a line and an isolated point.  Only the second
has the missing quadratic parameter class.  The local-algebra and
rank-locus checks are exact (\path{research/e80_singularity_recognition.py}).
\end{remark}

\section{The atlas}\label{sec:atlas}

Let $B=\am(\tau,k)$ be a continuous amplitude and $z=\sin^2B$; by the
first-time theorem $r(\cdot,k)$ has a unique inverse $\mathcal P(z,k)$ on
$(p_0,p_1)$, and
\begin{equation}\label{eq:sheet-param}
 \tau=F(B,k),\qquad \psi=\tau-\mathcal P(\sin^2B,k),\qquad t=2k\mathcal P(\sin^2B,k)
\end{equation}
parametrizes the first singular surface, one full period of $B$ covering
the four sign patches for one sign $s_2$.  The half-angle rectifying chart
of \cite{MoiseevSachkov2010Maxwell} is
\begin{equation}\label{eq:half-angle}
 W=\tan\frac\theta2,\qquad
 U=-\frac{y\cos\frac\theta2-x\sin\frac\theta2}{2\cos\frac\theta2},\qquad
 V=-\frac{x\cos\frac\theta2+y\sin\frac\theta2}{2\cos\frac\theta2},
\end{equation}
a diffeomorphism of $\{|\theta|<\pi\}$ onto $\mathbb R^3$.

\begin{proposition}[Stratified atlas]\label{prop:atlas}
With $q=\mathcal P-\epsi(\mathcal P)$, $s=\sn\mathcal P$, $d=\dn\mathcal P$,
$D=(1-k^2\sin^2B)^{1/2}$, the caustic for $s_2=+1$ is
\begin{equation}\label{eq:atlas-target}
 U=\frac{qD}{d},\qquad V=\frac{kf_1(\mathcal P)\sin B}{d},\qquad W=-\frac{ks\cos B}{d}.
\end{equation}
The other patches are $B\mapsto-B:(U,V,W)\mapsto(U,-V,W)$,
$B\mapsto\pi-B:(U,V,-W)$, $B\mapsto B+\pi:(U,-V,-W)$, and
$s_2\mapsto-s_2:(-U,-V,W)$.  On every open patch $(U,V)$ are local
coordinates, and on $0<B<\pi/2$
\begin{equation}\label{eq:orientation}
 \det\frac{\partial(U,V)}{\partial(z,k)}<0 .
\end{equation}
\end{proposition}
\begin{proof}
\eqref{eq:atlas-target} is \eqref{eq:exp} in the chart \eqref{eq:half-angle},
as an identity in $(p,\tau,k)$ with $p=\mathcal P$, $\sin B=\sn\tau$: with
$\psi=\tau-p$, $\psi_1=\tau+p$, $S=\sn\tau$ and $\Delta_0=1-mS^2s^2$, the
addition theorems give $\cn\psi-\cn\psi_1=2SsDd/\Delta_0$,
$\sn\psi=(Scd-CDs)/\Delta_0$, $\dn\psi=(Dd+mSsCc)/\Delta_0$, and the
addition theorem for $\epsi$ gives $\Delta=2q+2mS^2scd/\Delta_0$; with
$\sin\theta=k(\sn\psi\dn\psi_1-\dn\psi\sn\psi_1)$,
$\cos\theta=\dn\psi\dn\psi_1+m\sn\psi\sn\psi_1$ and
$W=\sin\theta/(1+\cos\theta)$, substitution into \eqref{eq:half-angle}
reduces, modulo $c^2+s^2=1$, $d^2=1-ms^2$ and the same relations at
$\tau$, to \eqref{eq:atlas-target} (\path{research/e79_atlas_identity.py},
symbolic; the formulas hold off the sheet as well).
The sign actions are Jacobi parity and periodicity.  The restriction of
$d\Exp$ to the tangent plane of $\Sigma$ is injective exactly when the
kernel $v$ is transverse to $\Sigma$, i.e. when $v(J_2)\ne0$: that is the
fold property, Theorem~\ref{thm:fold}, so the caustic is an immersed
surface.  Its projection determinant $\partial(U,V)/\partial(z,k)$ reduces
on $J_2=0$, after factors of fixed sign, to the same polynomial $P$ as the
certificate (Proposition~\ref{prop:app-projection} in the appendix), hence
is nonzero by Theorem~\ref{thm:fold}; so $(U,V)$ are local coordinates on
each open patch, and the sign in \eqref{eq:orientation} is fixed by
continuity and one evaluation (at $k=0.8$, $B=\pi/4$).
\end{proof}

\begin{figure}[t]
\centering
\includegraphics[width=\textwidth]{../figures/caustic_atlas.pdf}
\caption{The first caustic.  (a) The fundamental patch $0<B<\pi/2$,
$s_2=+1$, in the half-angle chart $(U,V)$ for seven moduli (color), with
the Maxwell seam $B=0$ (solid, the cusp edge, $V=0$) and the upper seam
$B=\pi/2$ (dashed) across all moduli, and the $D_4^+$ point at $k_0$.
(b) The planar shadow $(x,y)$ of the same curves and seams; each curve is
the projection of the first conjugate points of the geodesics of one
modulus.}
\label{fig:atlas}
\end{figure}

\section{Embedding and incidence}\label{sec:embedding}

\begin{lemma}[Slope]\label{lem:slope}
On the fundamental patch, write $\zeta=ab$ for the Riccati variable of
\eqref{eq:ab}, retaining $z=\sin^2B=r$ for the atlas coordinate.  Then
$H=V/U$ satisfies
\begin{equation}\label{eq:H}
 H=\frac k{k'}\,a\sqrt\zeta,\qquad
 \frac{\partial\log H}{\partial p}=\frac{a'}a+\frac{\zeta'}{2\zeta}>0.
\end{equation}
Since $z=r=\zeta/(1-m+m\zeta)$ increases with $\zeta$, also $H_z>0$.
At fixed $k$ the slope runs once from $0$ at the Maxwell
seam to an upper value $H_1(k)$ at the upper seam.
\end{lemma}
\begin{proof}
$V/U=kf_1S/(qD)$ and by \eqref{eq:r-in-z}
$S^2/D^2=r/(1-mr)=\zeta/(1-m)$.  The derivative is positive by
\eqref{eq:riccati} and \eqref{eq:a-below-g}; and $p$ increases with the
atlas coordinate $z=r$.
\end{proof}

\begin{lemma}[Upper slope]\label{lem:upper-slope}
$H_1(k)=k/k'$ for $k\le k_0$, and $H_1$ is strictly increasing on $(0,1)$,
with $H_1\to0$ as $k\downarrow0$ and $H_1\to\infty$ as $k\uparrow1$
(for $k>k_0$ the sign of $H_1'$ rests on an exact differentiation replayed
in \path{scripts/first_caustic_global_univalence_proof.py}).
\end{lemma}
\begin{proof}
For $k\le k_0$ the upper seam is $p=2K$, where $f_1=q=2(K-E)$ and $z=1$.
For $k>k_0$ it is $\alpha_1=0$, where $z=1$ and $H_1=(k/k')\,f_1/q$.  The
argument is as follows.  Set $\varpi=2K-p_\alpha$ and, at $\varpi$,
$\tilde s=\sn\varpi$, $\tilde g=\cn\varpi$, $\tilde d=\dn\varpi$,
$\tilde e=2E-\epsi(\varpi)$, so that $e=\tilde e$ and $q=p_\alpha-\tilde e$ at
$p_\alpha$; the endpoint equation $\alpha_1(p_\alpha)=0$ reads
$q(\tilde e\tilde s-\tilde g\tilde d)+\tilde d(\tilde g\tilde e+\tilde s\tilde d)=0$.
With $\upsilon=\tilde d/\tilde e$ and $\varrho=\tilde e\tilde s/(\tilde d\tilde g)\in(0,1)$,
\[
 \frac{f_1}q=\frac{\sqrt{1+\varrho^2\upsilon^2}}{1+\varrho\upsilon^2},\qquad
 \frac q{\tilde e}=\frac{1+\varrho\upsilon^2}{1-\varrho},\qquad
 \frac{p_\alpha}{\tilde e}=\frac{2-\varrho(1-\upsilon^2)}{1-\varrho},
\]
and exact differentiation along $\alpha_1=0$ shows that $H_1'(k)$ has the
sign of $-F_-G$ with
\[
 -F_-=1+\varrho\upsilon^2(1-\varrho)+\varrho^3\upsilon^4+k^2\varrho^3\upsilon^2(1-\upsilon^2)
 +2(1-k^2)\varrho^2\upsilon^2>0,
\]
\[
 G=(1-\varrho)(1+\varrho\upsilon^2)\{1+(2\varrho-1)\upsilon^2\}
 -(1-m)\varrho\upsilon^2\{2-\varrho(1-\upsilon^2)\}^2 .
\]
$G>0$ follows from $G>(1-\varrho)^2\mathcal G/\tilde e^4$ with
$\mathcal G=\tilde eq(\tilde e^2-\tilde d^2)-(1-m)\tilde d^2p_\alpha^2$, and from
the three bounds $k_0^2>3/4$, $\tilde e>3/2$, $2(1-m)(K-E)<3/4$ of
Lemma~\ref{lem:app-bounds} in the appendix, which, with $\tilde d\le1$,
$\tilde e<q<2(K-E)$ (so that $p_\alpha^2/q=q+2\tilde e+\tilde e^2/q<2(K-E)+3\tilde e$)
and $1-m<1/4$, give
$\mathcal G>q\bigl[\tilde e(\tilde e^2-1)-\tfrac34-\tfrac34\tilde e\bigr]
=\tfrac q4(\tilde e+1)(2\tilde e-3)(2\tilde e+1)>0$.
The two formulas for $H_1$ agree at $k_0$; the last factor of $f_1/q$
exceeds $1/2$, so $H_1\to\infty$; and $H_1\to0$ with $k$.
\end{proof}

\begin{theorem}[Each open patch is embedded]\label{thm:univalence}
On every open patch the projection $(z,k)\mapsto(U,V)$ is injective; hence
$(U,V,W)$ embeds the patch.
\end{theorem}
\begin{proof}
Take the fundamental patch.  For $h>0$ the level $H=h$ exists exactly for
$\kappa(h)<k<1$, $\kappa=H_1^{-1}$, by Lemma~\ref{lem:upper-slope} and is a single graph
$z=z_h(k)$ by Lemma~\ref{lem:slope}.  The quotient rule and
\eqref{eq:orientation} give
$\det\partial(U,H)/\partial(z,k)=U^{-1}\det\partial(U,V)/\partial(z,k)<0$,
so along the level
\[
 \frac{dU}{dk}\Big|_{H=h}=-\frac{\det\partial(U,H)/\partial(z,k)}{H_z}>0 .
\]
Two source points with the same $(U,V)$ have the same $H$ and $U$, hence
the same $k$, hence the same $z$.
\end{proof}

\begin{theorem}[Closed incidence]\label{thm:closed}
For $0<k<1$, the caustic is a topological embedding of every fundamental
strip $0\le B\le\pi/2$, including its seams, into $\SE$.  This embedding
is proper into the punctured target $\SE\setminus\{e_{\SE}\}$, where
$e_{\SE}$ is the identity, but not into $\SE$ itself.  The Maxwell seam
is cuspidal; the upper seam is cuspidal for $k<k_0$, an umbilic at $k_0$,
and a fold for $k>k_0$ (Theorem~\ref{thm:edges}, Proposition~\ref{prop:d4}).  The eight open patches and their seams have no target
incidences except the identification of a common source seam.  As
$k\downarrow0$ the caustic collapses to the identity; no sequence with
$k\uparrow1$ stays in a compact set.
\end{theorem}
\begin{proof}
On the lower seam $B=0$ one has $f_1=0$, so $U=q/d=-s/c$, $V=0$,
$-W=ks/d=kU/\sqrt{1+(1-m)U^2}$, and along $f_1=0$
\[
 U_k(f_1)_A-U_A(f_1)_k=\frac{k\beta_1}{c^2(k^2-1)}>0,\qquad (f_1)_A=-\frac dc>0,
\]
($A=\am p$), so $dU/dk>0$; $U\to0$ as $k\downarrow0$ and
$U\ge q>K-2E\to\infty$ as $k\uparrow1$: the seam is embedded.  The upper
seam is embedded because its slope $H_1$ is strictly increasing.  The signs
on the three strata are $(U,V,W)=(+,+,-)$ inside, $(+,0,-)$ on $B=0$,
$(+,+,0)$ on $B=\pi/2$; the eight open patches occupy the eight open
octants, and adjacent patches share seams: the four lower seams carry the
labels $(\pm,0,\pm)$ and the four upper seams $(\pm,\pm,0)$, one seam per
label.  As $k\downarrow0$,
uniformly on the strip, $p=\pi+O(k^2)$, $q=O(k^2)$, $\sn p=O(k^2)$, so
$(U,V,W)=O(k^2,k^3,k^3)$ and $t=2kp\to0$.  For $k\uparrow1$ let
$R^2=x^2+y^2=4(U^2+V^2)/(1+W^2)$ be the spatial radius; from
\eqref{eq:atlas-target} and \eqref{eq:r-in-z}, with the Riccati variable
$\zeta=ab$, $\chi=k^2g^2\zeta/(1-m)$ and $\varrho_0=a/g$,
\[
 \frac{R^2}{4q^2}=1-\frac{(1-\varrho_0^2)\chi}{1+\chi}\ge1-\frac2q,
\]
because $1-\varrho_0=sd/(qg)\le q^{-1}\sqrt{1+1/\chi}$, i.e.
$s^2d^2\le g^2+(1-m)/(m\zeta)$, which follows from
$s^2d^2=(1-y)(1-m+my)\le1-m+my\le y+(1-m)/m\le y+(1-m)/(m\zeta)$ with
$y=g^2$, and $\sqrt{\chi(1+\chi)}\le1+\chi$; and $q\ge q(p_0)>K-2E\to\infty$.
The two boundary limits confine the inverse image of any compact subset
of $\SE\setminus\{e_{\SE}\}$ to
$[0,\pi/2]\times[\epsilon,1-\epsilon]$ for some $\epsilon>0$.
It is closed there, hence compact.  Properness and the proved
injectivity give a topological embedding into the punctured target,
and therefore into $\SE$.  Properness into $\SE$ fails: at fixed $B$,
a sequence $k_j\downarrow0$ has images converging to $e_{\SE}$, so
these images together with $e_{\SE}$ form a compact set whose inverse
image contains a sequence without a convergent subsequence in the strip.
\end{proof}

\section{Every sheet}\label{sec:every-sheet}

By \cite[Theorem 4.1]{SE2FirstConjugate}, in the cell $[p_n^1,p_{n+1}^1)$,
$n\ge1$, every rotating geodesic has one conjugate time on each of the brackets
(two distinct times except at $k=k_0$, $\psi\equiv K\pmod{2K}$,
where both equal $2nK$)
$B_n^-=[p_n^1,\min(2nK,p_n^\alpha)]$ and $B_n^+=[\max(2nK,p_n^\alpha),b_n]$,
where $p_n^\alpha$ is the zero of $\alpha_1$ in the cell and $b_n$ the
zero of $\beta_1$ in $(2nK,(2n+1)K)$; on each open bracket $0<r<1$ and
the conjugate points are $\sn^2\tau=r$.  The \emph{sheets} $\Sigma_n^\pm$
are these two surfaces in the source, each with two mirror branches
$\tau=\pm u$; their closures add the seams $S=0$ (at $p_n^1$ for
$\Sigma_n^-$, at $b_n$ for $\Sigma_n^+$) and $C=0$ (at $2nK$ or $p_n^\alpha$).
The first sheet is $\Sigma_1^-$.  With the sign
$\sigma_n=(-1)^{n+1}$ and $g=|\cn p|$, $\mathsf u=|\sn p|\dn p/q$, $\mathsf v=|\sn p|\dn p/e$,
\begin{equation}\label{eq:sheet-signs}
 (\sigma_na,\sigma_nb)=(g-\mathsf u,\ g+\mathsf v)\ \text{on }\Sigma_n^-,\qquad
 (\sigma_na,\sigma_nb)=(g+\mathsf u,\ g-\mathsf v)\ \text{on }\Sigma_n^+,\qquad \sigma_nc=-g,
\end{equation}
because $f_1$ has the sign $\sigma_n$ on the cell, $\beta_1$ the sign
$(-1)^n=-\sigma_n$ on $(b_{n-1},b_n)$, and $\sn p$ the sign $\sigma_n$ on
the left half and $-\sigma_n$ on the right half of the cell.  Figure~\ref{fig:sheets}
shows the first three cells.

\begin{theorem}[Every sheet]\label{thm:every-sheet}
Let $n\ge1$ and $0<k<1$ (at $k=k_0$ the two brackets touch at $2nK$ and
$0<r<1$ still holds on both open brackets,
\cite[Theorem 4.1]{SE2FirstConjugate}).
\begin{enumerate}[nosep,leftmargin=*]
\item[(i)] On the open sheets $\Sigma_n^\pm$ the map $\Exp$ has corank one,
      its kernel is $(\theta_k,-\theta_\psi,0)$, and the identity
      \eqref{eq:fold-identity} holds with $S=\sn\tau$ and $s=\sn p$ signed.
\item[(ii)] (computer-assisted) $Q>0$ on $\Sigma_n^-$ and $Q<0$ on $\Sigma_n^+$, up to and
      including their $\alpha_1$ ends when $k\ne k_0$.  Hence every open
      sheet is a fold surface, and every noncritical $\alpha_1$ seam is a fold.
      At $k=k_0$ the common end has $Q=0$ and is covered by (iii).
\item[(iii)] The seam $S=0$ of $\Sigma_n^-$ (at $p_n^1$) is an $A_2$ cusp
      edge; the seam $S=0$ of $\Sigma_n^+$ (at $b_n$) is a fold; the seam
      $C=0$ at $p=2nK$ is an $A_2$ cusp edge for $k\ne k_0$; at
      $(p,\tau,k)=(2nK,K,k_0)$ the map has corank two and is a hyperbolic
      umbilic $D_4^+$.
\end{enumerate}
\end{theorem}
\begin{proof}
(i) Proposition~\ref{prop:kernel} and Theorem~\ref{thm:fold-identity} used
only $J_2=0$, $\sn p\ne0$ and $\sn\tau\ne0$, which hold on every open
sheet ($p\notin K\mathbb Z$ there, and $0<r<1$).

(ii) In the variables \eqref{eq:sheet-signs} the certificate of
Theorem~\ref{thm:fold} applies to $\Sigma_n^-$ verbatim: $0<\mathsf u<g<1$
and $\mathsf v>0$ by \eqref{eq:sheet-signs}; $ab\le1$ is $\mathsf Z\ge0$
(equality at the $\alpha_1$ end); $e=\epsi(p)>E>1$ because $p>K$; and
$me-(1-m)q=m\int_0^p\cn^2>0$.  So $P=S_0+\mathsf B S_B+\mathsf Z S_Z>0$ and
$Q>0$, with $\mathsf Z=0$ allowed at the end.  On $\Sigma_n^+$ the
orientation of \eqref{eq:sheet-signs} is reversed, and $Q<0$ there, up to
and including the $\alpha_1$ end when $k<k_0$, by Proposition~\ref{prop:app-right} in
the appendix: the bound $\mathsf t=L_0\mathsf v/\mathsf u<1/2$ on every
later cell, the concavity of the scaled $-Q$ in $g^2$, an elementary
argument on the face $ab=1$ and four exact rational certificates on the
face $\mathsf v=\rho/2$ (\path{data/universal_right_sv_certificates.json},
replayed by \path{scripts/universal_right_fold_proof.py}).  The hypotheses of both
certificates were also checked on $4{,}480$ interior points of the brackets
of cells $1$--$8$ at seven moduli.  The minima of the five recorded
diagnostic margins lie between $5\times10^{-4}$ and $2\times10^{-2}$
(\path{research/e75_step3_certificate_domains.py}); this range does not
bound every polynomial margin.  All points were independently replayed
at $60$ digits (\path{research/e80_certificate_domain_audit.py}).
This finite survey does not test the critical endpoint.  Away from
$k=k_0$, the $\alpha_1$ seams are interior points of the smooth singular
surface where the two mirror branches meet ($C=0$, $S=1$), so (i) applies
and they are folds by \eqref{eq:fold-identity}.  At the critical end,
$s=0$, $e=q$, $a=b=-c=\pm1$ instead give $Q=0$.

(iii) At $S=0$ the formulas of Theorem~\ref{thm:edges}(i) hold for every
$n$: $x_\psi=-2kCdf_1$, $y_\psi=2msf_1$, $\theta_\psi=0$,
$\theta_k=2C\beta_1/(1-m)$.  At $p_n^1$ therefore $\eta_\psi=0$, the kernel
is the phase direction, $Q=-3(1-m)bc\ne0$, and
$v^2(J_2)=-12(1-m)q^4e^3\beta_1^2\,c/(s\,d\,\alpha^2)$, nonzero because
$c/s<0$ on $((2n-1)K,2nK)$: a cusp, with the rank condition as before.
At $b_n$ one has $\theta_\psi=\theta_k=0$ but $x_\psi\ne0$, so the
$x$-row extension $v=(x_k,-x_\psi,0)$ is the kernel, and with
$\beta_1(b_n(k),k)=0$ an exact computation in the Jacobi algebra
(\path{research/e74_step3_edges.py}, extended to $b_n$ by quasi-periodicity) gives
\begin{equation}\label{eq:beta-seam}
 v(J_2)=2kC^3d\,f_1^2\,\partial_k^{\rm src}\beta_1
 =\frac{2k^2C^3d\,f_1^2}{mc}\,(es^2-qc^2),
\end{equation}
where $\partial^{\rm src}_k$ is the derivative at fixed $(\psi,t)$.  With
the signs normalized so that $s,c,d>0$, the seam equation $ce=ds$ gives
$c^2=d^2/(e^2+d^2)$ and $d^2(d^2+e^2-1)=(1-m)e^2$, hence
$es^2-qc^2=(e^3-qd^2)/(e^2+d^2)$, and $qd^2<e^3$ follows from
$(1-m)q<(1-m)(2n+1)K\le(2n+1)\pi/2<2n(4n^2-1)<e(e^2-1)$, using
$e>2nE>2n$: the seam is a fold.  At $(2nK,K)$, $\eta_\psi=0$, $d\theta=0$
and $\eta_k=2(e-q)\bigl((-1)^{n+1},-k/k',0\bigr)$ with $e=2nE$,
$q=2(K-E)n$, so $e-q=2n(2E-K)$ vanishes exactly at $k_0$; for $k\ne k_0$
the computation of Theorem~\ref{thm:edges}(ii) goes through with the sign
$(-1)^{n+1}$ cancelling in $v^2(J_2)=-12eq(1-m)(e-q)^2$: a cusp.  At $k_0$
the generating family of Proposition~\ref{prop:d4} applies with
$e=\epsi(2nK)=nK>1$ (for even $n$ after the target reflection
$(V,W)\mapsto(-V,-W)$), its discriminant $-e^4k^8(e^2-1)/(12k'^2)$ and
versal coefficient $k/[2k'(e^2-1)]$ keeping their signs; the intrinsic
test of Remark~\ref{rem:intrinsic} confirms $D_4^+$ at $n=1,2,3$; the
closed forms hold for every $n$ because at $p=2nK$, $\tau=K$ the Jacobi
values are those at $p=0$ up to the sign $(-1)^n$ (absorbed by the
reflection for even $n$), and $q=e=nK$ at $k_0$, so the contractions
depend on $(e,k)$ only.
\end{proof}

\begin{figure}[t]
\centering
\includegraphics[width=\textwidth]{../figures/caustic_sheets.pdf}
\caption{The sheets of the first three cells at $k=0.8$, both mirror
branches: (a) in the half-angle chart, solid for $B_n^-$ and dashed for
$B_n^+$; (b) their planar shadow.  Each bracket sheet is a fold surface;
its $S=0$ seam is a cusp edge for $B_n^-$ and a fold for $B_n^+$, and its
$C=0$ seam is the cusp edge at $2nK$ or the fold at $p_n^\alpha$
(Theorem~\ref{thm:every-sheet}).}
\label{fig:sheets}
\end{figure}

\begin{proposition}[Leading term of $Q$]\label{prop:leading}
Put $g=-\sigma_nc$ and $\mathsf u=g-\sigma_na$, $\mathsf v=\sigma_nb-g$, so that
$\mathsf u=\rho/q$, $\mathsf v=\rho/e$ with $\rho=\sigma_n\sn p\dn p$.  Then $Q$
is a polynomial in $(\mathsf u,\mathsf v)$ of degree at most three in each
variable, with twelve terms and coefficients in $(g,m)$, vanishing at
$\mathsf u=\mathsf v=0$, with linear part
$4g(1-g^2)(1-m+mg^2)\,\mathsf u$ and no linear term in $\mathsf v$.  On the
cell $n$ at $p=2nK+\sigma$, $|\sigma|<K$, exactly
$q=2n(K-E)+\sigma-\epsi(\sigma)$ and $e=2nE+\epsi(\sigma)$, hence
\begin{equation}\label{eq:Q-leading}
 Q=\frac{2\,\sn^3p\,\dn p\,(-\cn p)\,(1-m+m\cn^2p)}{n\,(K-E)}+O(n^{-2}),
\end{equation}
uniformly for $\sigma$ in compact subsets of $(-K,0)\cup(0,K)$.  Its
sign is that of $-\sn\sigma\cn\sigma$: positive on the left half of the
cell, negative on the right half, at every modulus.
\end{proposition}
\begin{proof}
Substituting $a=\sigma_n(g-\mathsf u)$, $b=\sigma_n(g+\mathsf v)$, $c=-\sigma_ng$
into \eqref{eq:Q} and reducing with $s^2=1-g^2$ gives the twelve-term
polynomial; its constant term vanishes and its linear part is as stated
(\path{research/e80_late_polynomial.py}, exact).  Quasi-periodicity gives the exact
$q$ and $e$, so $\mathsf u,\mathsf v=O(1/n)$ and the linear term dominates.
\end{proof}

\begin{corollary}[Late conjugate points are folds]\label{cor:late-folds}
The conjugate points of every cell $n\ge1$ are folds, for every
$0<k<1$ and every phase off the seams.  On every compact set of $(\psi,k)$
with $\psi\not\equiv K$ there is $n_0$ such that for $n\ge n_0$ this follows
from \eqref{eq:Q-leading} alone, without the certificates.
\end{corollary}
\begin{proof}
The first sentence is Theorem~\ref{thm:every-sheet}(ii).  By the
large-index law \cite[Theorem 5.1(ii)]{SE2FirstConjugate} the offsets
$\sigma$ of the two conjugate points of cell $n$ converge to
$\mp K/2-\psi/2$, with $\psi$ taken in $(-K,K)$, uniformly on the set
(the remainders of the expansions behind that law are explicit continuous
functions of $(\sigma,k)$), so for $n\ge n_0$ they lie in a compact subset
of $(-K,0)\cup(0,K)$, where \eqref{eq:Q-leading} has the sign of
$-\sn\sigma\cn\sigma$ and its remainder is smaller.
\end{proof}

\begin{remark}
For $k\ne k_0$ the sign of $Q$ changes between the two brackets inside
the gap between $2nK$ and $p_n^\alpha$, in either order.  Put
$D_0=K-E$ and $\sigma=w/n$.  The leading polynomial balance is
\[
 n^4Q\approx w^2\left[-\frac{2w}{D_0}
       +\frac{3(2E-K)}{4D_0^2E}\right],\qquad
 n\sigma_\alpha\approx\frac{2E-K}{2D_0E}.
\]
Its nonzero root is three quarters of the endpoint offset.  The
numerically computed argument-offset ratios
$(p_Q-2nK)/(p_n^\alpha-2nK)$ are $0.771$, $0.755$, $0.7503$
at $n=2,5,20$ for $k=0.6$; these are not ratios of Jacobi $\sn$ values.
The balance is checked in \path{research/e80_late_polynomial.py};
we do not use it as a certificate in the shrinking gap, since the
fold certificates already cover every open sheet.
The two cusp seams lie on the closures of the reflection Maxwell strata
$\operatorname{MAX}_2$ ($p=p_j^1$) and $\operatorname{MAX}_6$ ($p=2jK$) of
\cite[Theorem 6.2]{MoiseevSachkov2010Maxwell}, the fold seams on neither.
\end{remark}

\section{Incidences between sheets}\label{sec:incidence}

Two sheets meet where a point of $\SE$ is a conjugate point along two
different geodesics.  The chart \eqref{eq:half-angle} first restricts
which sheets can meet at all, and the scalar inverse of the endpoint map
\cite{SE2Counts} then reduces the question to the level function of a
target.

\paragraph{Notation.}  In this section $(A_q,B_q)$ are the coordinates of
the target in the frame of the heading bisector (the amplitudes $A=\am p$,
$B=\am\tau$ of Sections~\ref{sec:atlas}--\ref{sec:embedding} do not occur
here), $\nu$ is the parameter of the chart curve (the kernel $v$ of
Section~\ref{sec:kernel} does not occur), $L(m)=4(K(m)-E(m))$ is the drift
per period of \cite{SE2Counts}, written $L$ when the argument is clear, and
$\beta_N=4E(m)/L(m)$; $q$ denotes a target, $\ell$ a real level, and the
drift $q=p-\epsi(p)$ and $\ell=e\,(p-\epsi(p))$ of \eqref{eq:ab} do not occur;
$s_i,c_i,d_i$ are the scaled values defined below, not those of
Section~\ref{sec:fold}.

Let $F(n,\pm)$ be the \emph{fundamental patch} of $\Sigma_n^\pm$: the
branch $\tau=u\in(0,K)$, $s_2=+1$.  The group $G$ of order eight generated
by $(U,V,W)\mapsto(U,V,-W)$ (the branch $\tau=2K-u$), $(U,-V,W)$
($\tau=4K-u$) and $(-U,-V,W)$ ($s_2=-1$) carries the fundamental patches
onto the whole conjugate locus.

\begin{lemma}[Octants]\label{lem:octants}
On every open sheet $U>0$.  On $F(n,\pm)$, $V$ has the sign
$\sigma_n=(-1)^{n+1}$ and $W$ the sign $-\sigma_n$ on $F(n,-)$, $+\sigma_n$
on $F(n,+)$.  Hence each fundamental patch lies in one open octant, the
$s_2=-1$ copies lie in $U<0$ and never meet the $s_2=+1$ ones, and
$F(n,b)$ can meet $g\,F(m,b')$ only for the single $g=(s_V,s_W)\in G$ with
\[
 s_V=\sigma_n\sigma_m,\qquad s_W=\epsilon_b\epsilon_{b'}\sigma_n\sigma_m,
 \qquad \epsilon_-=-1,\ \epsilon_+=+1 .
\]
\end{lemma}
\begin{proof}
$U=qD/d>0$ since $q,D,d>0$; $V=kf_1S/d$ with $\operatorname{sgn}f_1=\sigma_n$
on the cell and $S=\sn u>0$; $W=-ksC/d$ with $C=\cn u>0$ and
$\operatorname{sgn}\sn p=\sigma_n$ on the left half of the cell, $-\sigma_n$
on the right half.  The seams $S=0$, $C=0$ have $V=0$, $W=0$ and are not in
the open patches.
\end{proof}

\begin{proposition}[Conjugate preimages and the level function]\label{prop:level}
Let $q=(X,\theta)$, $X\ne0$, $0<|\theta|<\pi$, be off the seam of the
scalar inverse \cite[Proposition 2.1]{SE2Counts}.  Its rotating preimages
with $s_2=+1$ are the points $z$ of its admissible intervals at which the
level function
\begin{equation}\label{eq:level-function}
 N(z)=\frac{A_q-Q_0}{L(m)},\qquad L(m)=4\bigl(K(m)-E(m)\bigr),
\end{equation}
is a nonnegative integer \cite{SE2Counts}; here $A_q>0$ is the drift of
the target along the axis of its elastica, $L(m)$ the drift per period,
and $Q_0\in(0,L(m))$ the drift of the residual argument
$\delta(z)\in(0,4K)$ (Proposition~\ref{prop:N-prime} gives $Q_0$ in closed
form).  At such a point the source has phase $\psi(z)$, modulus $k(z)$,
residual argument $\delta(z)$ and time $t=k\delta+N\cdot4kK$.  Then:
\begin{enumerate}[nosep,leftmargin=*]
\item[(i)] the preimage is a conjugate point if and only if $N'(z)=0$
      \cite[Proposition 2.2]{SE2Counts};
\item[(ii)] a conjugate preimage on $\Sigma_n^-$ has $N=n-1$ and
      $\delta>2K$; one on $\Sigma_n^+$ has $N=n$ and $\delta<2K$;
\item[(iii)] $q$ lies on the image of $\Sigma_n^-$ if and only if $N$ has a
      critical point at level $n-1$ with $\delta>2K$, on the image of
      $\Sigma_j^+$ if and only if $N$ has one at level $j$ with $\delta<2K$,
      and on both if and only if both occur.
\end{enumerate}
\end{proposition}
\begin{proof}
(i) is the regularity bridge of the counts paper.  (ii) $p=t/(2k)=2NK+\delta/2$,
so $p/(2K)=N+\delta/(4K)$; on the left bracket $p\in((2n-1)K,2nK)$, whence
$N=n-1$ and $\delta\in(2K,4K)$; on the right bracket $p\in(2nK,(2n+1)K)$,
whence $N=n$ and $\delta\in(0,2K)$.  (iii) restates (i) and (ii).
\end{proof}

\paragraph{The incidence map.}  For cells $1$--$4$ the eight fundamental
patches were tested against the $31$ other patches of these cells (four
branches of each bracket, $s_2=+1$): on $25$ slices $k_1=\mathrm{const}$
of the first patch, up to four nearest sampled approaches within distance
$0.35$ were refined by least squares on $(f_1,f_2,k_2)$, where $f_i$ is
the fractional position in its bracket (with interpolated bracket
endpoints), three equations in three unknowns, and
accepted at residual $10^{-10}$ with interior parameters
(\path{research/e76_incidence_trace.py}).  The result
(Figure~\ref{fig:incidence}): among the $248$ ordered pairs the search
detected intersections for the following ten unordered pairs
\[
 F(n,-)\ \cap\ g\,F(j,+)\ne\emptyset,\qquad 1\le n\le j\le4,
\]
each detected from both sides, with $g$ the reflection
$W\mapsto-W$ for $n\equiv j$ and $V\mapsto-V$ otherwise, as
Lemma~\ref{lem:octants} dictates.  No left--left or right--right
intersection, pair with $j<n$, or self-intersection was detected.
The sampled same-cell traces are consistent with branches from the
umbilic toward $k=1$; this finite search does not establish their global
endpoints or exhaust all intersections.  At the sampled intersections
the left sheet has the larger modulus, and the right-sheet modulus
decreases with $j-n$ (from about $0.9$ at $j=n$ to about $0.3$ at $(1,4)$).
A replay of all $123$ saved accepted points recovers target residuals
below $6.1\times10^{-13}$.  The finite-grid detector
(\path{research/e76_level_function.py}) finds exactly two critical points,
both maxima, at $121$ targets; at one further target it finds two maxima
and two minima below level zero, and at the remaining target it finds
only one maximum.  Independent $60$-digit differentiation confirms the
negative minima and recovers the missed maximum near level $1$ at the
last target (the detected maximum is near level $2$).  These tests do
not certify completeness of the critical-point search.

\begin{figure}[!htb]
\centering
\includegraphics[width=\textwidth]{../figures/caustic_incidence_curves.pdf}
\caption{Incidences between the sheets of cells $1$--$4$.  (a) The
accepted intersection points in the planar shadow, over the sheets in
grey; (b) the same points in the $(k_1,k_2)$ plane of the moduli of the
two patches, in the order of the label (\texttt{patch A}, \texttt{patch B}
on the axes); each unordered pair appears twice, mirrored in the
diagonal, and each color serves two pairs (see the legend).  The sampled
traces are monotone, and the trace with the left sheet first lies below
the diagonal: the left sheet has the larger modulus at these points.
Labels: \texttt{n1-id x n2+V} is
$F(1,-)\cap g F(2,+)$; the token after each cell index and bracket sign
names the element of $G$ applied to that patch, \texttt{id} the identity
(the fundamental patch itself), \texttt{V} the reflection $V\mapsto-V$
and \texttt{W} the reflection $W\mapsto-W$.}
\label{fig:incidence}
\end{figure}

Two facts about the level function are theorems.  Work in the frame of
the heading bisector: let $\phi$ be the angle of the elastica axis from
the bisector, $\omega=2\phi$, let $(A_q,B_q)$ be the along-axis and
transverse coordinates of the target, $R^2=A_q^2+B_q^2$, and let $\nu$ be
the polar angle of the target seen from the axis, $A_q=R\sin(\nu/2)$,
$B_q=R\cos(\nu/2)$; along the chart curve of a fixed target $\nu$ is the
parameter.  Write $s_i=k\sn$, $c_i=k\cn$,
$d_i=\dn$ at the initial ($i=0$) and final ($i=1$) phase, so
that $s_0=\sin a$, $s_1=\sin(a-\theta)$ for the axis angle $a$, and
$B_q=c_0-c_1$.

\begin{lemma}[The modulus along the chart curve]\label{lem:modulus}
Along the chart curve
\[
 4m=B_q^2+2-2\cos\theta\cos\omega+\frac{\sin^2\theta\,\sin^2\omega}{B_q^2},
\]
the chart is admissible only if $B_q^2$ lies strictly between
$(1-\cos\theta)(1-\cos\omega)$ and $(1+\cos\theta)(1+\cos\omega)$ (this is
$0<m<1$; admissibility also requires the sign conditions of
\cite[Section~2]{SE2Counts}), and
\begin{equation}\label{eq:m-prime}
\begin{aligned}
 B_q\,\partial_\nu m&=A_q\,c_0c_1+c_0s_1d_1-c_1s_0d_0,\\
 B_q\,\partial_\nu^2 m&=\frac{m(1-m)\,(c_0^3-c_1^3)}{2\,c_0^2c_1^2}
 \quad\text{at every critical point of } m .
\end{aligned}
\end{equation}
Hence every critical point of $m$ with $0<m<1$ is a strict minimum, and
$m$ is unimodal on every admissible interval.
\end{lemma}
\begin{proof}
The transverse coordinate is $B_q=c_0-c_1$ and both endpoints lie on the
circle $s_i^2+c_i^2=m$ with $s_0^2-s_1^2=\sin\theta\sin\omega$, which
gives $c_0=(B_q-\sin\theta\sin\omega/B_q)/2$ and the closed form; $m=1$ is
$B_q^2=4\cos^2(\theta/2)\cos^2\phi$ or $B_q^2=4\sin^2(\theta/2)\sin^2\phi$.
Along the curve $\partial_\nu A_q=B_q/2$, $\partial_\nu B_q=-A_q/2$,
$\partial_\nu s_i=d_i/2$, $\partial_\nu d_i=-s_i/2$, and
$\partial_\nu c_0=(s_1d_1-s_0d_0+A_qc_1)/(2B_q)$,
$\partial_\nu c_1=\partial_\nu c_0+A_q/2$ from $s_0^2+c_0^2=s_1^2+c_1^2$.
Differentiating the closed form gives the first identity
in~\eqref{eq:m-prime}.  Differentiating it once more and substituting
the critical value $A_q=s_0d_0/c_0-s_1d_1/c_1$ gives
$2B_q\,\partial_\nu^2m=c_0c_1(c_0-c_1)+c_0(1-2s_1^2+s_1^2d_1^2/c_1^2)
-c_1(1-2s_0^2+s_0^2d_0^2/c_0^2)$; the identity
$(1-2s^2)c^2+s^2d^2=c^4+m(1-m)$ on the circle turns the brackets into
$c_i^2+m(1-m)/c_i^2$, the cubic terms cancel, and the second identity
follows.  Its sign is that of $B_q$, so $\partial_\nu^2m>0$ at every
critical point with $0<m<1$; two strict minima on an open interval would
enclose a maximum.
\end{proof}

\begin{proposition}[Derivative of the level function]\label{prop:N-prime}
Let $w_i=s_id_i/c_i$, $A_q^*=w_0-w_1$, and
$\Delta E=E(\varphi_1|m)-E(\varphi_0|m)$ the increment of the incomplete
integral of the second kind between the endpoint amplitudes, so that
$Q_0=\delta-\Delta E$.  On the chart domain
$J=(-\pi/2,\pi/2)\cap(\theta-\pi/2,\theta+\pi/2)$ of the principal axis
angle of Lemma~\ref{lem:positive-superlevels}, where $d_0,d_1>0$, along
the chart curve
\begin{equation}\label{eq:N-prime}
 L\,\partial_\nu N=-\frac{mB_q}{2c_0c_1}
 -\frac{\partial_\nu m}{2(1-m)}\Bigl(A_q^*+\Delta E+4N\,E(m)\Bigr).
\end{equation}
Consequently $\partial_\nu N\ne0$ at the minimum of $m$, with the sign of
$-B_qc_0c_1$; the critical points of $N$ are the roots $\Xi=A_q-A_q^*$ of
\[
 \beta_N \Xi^2+\bigl[A_q^*+\Delta E+\beta_N(A_q^*-Q_0)\bigr]\Xi+\frac{m(1-m)B_q^2}{c_0^2c_1^2}=0,
 \qquad \beta_N=\frac{4E(m)}{L(m)}.
\]
\end{proposition}
\begin{proof}
$Q_0=\mathcal D(\varphi_1|m)-\mathcal D(\varphi_0|m)$ with $\mathcal D=F-E$,
$w(\varphi)=\sqrt{1-m\sin^2\varphi}$, $\partial_\varphi
\mathcal D=m\sin^2\varphi/w$, $\partial_m\mathcal D=[E(\varphi|m)-\sin\varphi\cos\varphi/w]/(2(1-m))$,
and the amplitude derivatives $\partial_\nu\varphi_i=(md_i-s_i\partial_\nu m)/(2mc_i)$.
Collecting, the algebraic part of $\partial_\nu A_q-\partial_\nu Q_0$ is
$B_q/2-s_1^2/(2c_1)+s_0^2/(2c_0)=-mB_q/(2c_0c_1)$ and the rest is
$\partial_\nu m$ times $-[A_q^*+\Delta E]/(2(1-m))$, using
$s_i^3/(c_id_i)/(2m)+s_ic_i/(2md_i(1-m))=w_i/(2(1-m))$; the term
$-N\partial_\nu L$ with $\partial_mL=2E/(1-m)$ completes~\eqref{eq:N-prime}.
Multiplying $\partial_\nu N=0$ by $\Xi$ and eliminating $N=(\Xi+A_q^*-Q_0)/L$
gives the quadratic.  The identities are checked
numerically in \path{research/e77_h1_proof_check.py} and
\path{research/e77_Nprime_formula.py}.
\end{proof}

The quadratic makes the critical points of $N$ a double cover of the
three-dimensional endpoint data $(m,\varphi_0,\varphi_1)$, so the global
structure of the level function can be scanned exhaustively.  The scan
shows that unimodality at every real level fails near zero heading, and
the conjecture is stated in the form the counts need.

\subsection*{A conditional count criterion and its certificate}

The critical points of the level function form a double cover of the
endpoint data $(m,\varphi_0,\varphi_1)$ (Proposition~\ref{prop:N-prime}),
so the sign of $N''$ at every critical point is a statement about an
explicit function of three variables.  The proposed interval verification
requires a complete replay: the previous run used incorrect incomplete-integral
enclosures at negative amplitudes and incorrect Taylor tightening of
quadratic roots near coalescence.  The repaired primitives pass independent
checks, but those checks do not certify the full domain.  Write $\mu=(\varphi_0+\varphi_1)/2$ and
$\Delta\varphi=\varphi_1-\varphi_0$; the statement is invariant under
$\mu\mapsto\mu+\pi$ and $\mu\mapsto-\mu$, so $\mu\in[0,\pi/2]$ is a
fundamental domain.  Let $\mathcal B$ be the box
\[
 \mathcal B=\{\,0.02\le m\le 0.98,\ 0\le\mu\le\pi/2,\ 0.02\le\Delta\varphi\le 2\pi-0.02\,\}
\]
with five sets removed: the slab $\mathcal S=\{|B_q|<0.02\}$ around the
bisector boundary $B_q=0$ (the seam of the scalar inverse, reached as
$\mu\to0$ and as $\Delta\varphi\to0,2\pi$), the two corners
$\mathcal C_1=\{\mu<0.08,\ \Delta\varphi<0.2\}$ and
$\mathcal C_2=\{\mu<0.08,\ \Delta\varphi>2\pi-0.2\}$ near short or
full-period residual pieces, and the tube
$\mathcal C_3=\{\mu<0.08,\ |\Delta\varphi-\pi|<0.1\}$ around the double
flip $c_0=c_1=0$, where both roots of the critical quadratic escape to
infinity, and the low-modulus set
\[
 \mathcal C_4=\{0.02\le m\le0.06,\ 0.94\le\mu\le\pi/2,\ 0.41\le\Delta\varphi\le0.80\}
 \ \cup\ \{0.78\le m\le0.82,\ 1.49\le\mu\le\pi/2,\ 0.02\le\Delta\varphi\le0.10\},
\]
the union of the $32$ grid cells (a wedge of $31$ cells and one further
cell) in which the budget passes of the interval engine leave boxes
open.  On the wedge the small critical root is an $O(m)$ difference of
$O(1)$ quantities and the interval enclosures widen faster than the
margin of \eqref{eq:verified-predicate} shrinks; a $50$-digit pointwise
survey finds $mN''(\nu_*)\to-c(\mu,\Delta\varphi)$ with $c\ge0.0013$ on the
wedge, so the first alternative holds there numerically with a margin
growing like $1/m$, but this is not certified and $\mathcal C_4$ is
declared as an exclusion.  By \emph{the family} we mean below the grid of
$100\times120\times240$ endpoint data $(m,\mu,\Delta\varphi)$ of
\path{research/e78_statement_L.py}, which samples the ambient endpoint
domain, including the removed sets, corners and tube; on this finite grid the tube
carries no interior minimum, its critical points having $R\ge0.4$ and
$N\ge0.07$.

\paragraph{Certificate.}
The computation establishes, at every interior critical
point $\nu_*$ with $A_q>0$ and endpoint data in
$\mathcal B\setminus(\mathcal S\cup\mathcal C_1\cup\mathcal C_2\cup\mathcal C_3\cup\mathcal C_4)$,
the predicate
\begin{equation}\label{eq:verified-predicate}
 N''(\nu_*)<0\qquad\text{or}\qquad
 N(\nu_*)<\tfrac12\ \ \text{and}\ \ R<L\bigl(m(\nu_*)\bigr),
\end{equation}
where $R=\sqrt{A_q^2+B_q^2}$.  The run of 21 September 2026 described in
the numerical paragraph below, with the repaired primitives, closes every
one of the $38{,}400$ cells outside the five sets: the $304$ boxes left
open by its budget passes (total volume $0.049\%$ of $\mathcal B$) all lie
in $\mathcal C_4$, $285$ in the wedge and $19$ in the further cell, and
are listed in \path{data/e80_full_grid_open_cells.json}.  The excluded
sets occupy $2.56\%$ of $\mathcal B$ (slab $2.06\%$, corners $0.25\%$,
tube $0.14\%$, $\mathcal C_4$ $0.11\%$).
The corollary below is an analytic implication of the predicate; off
$\mathcal S\cup\mathcal C_1\cup\mathcal C_2\cup\mathcal C_3\cup\mathcal C_4$
its hypothesis is the certified statement.

\begin{lemma}[Compactness of positive superlevels; {\cite[Proposition~2.4]{SE2Counts}}]\label{lem:positive-superlevels}
Fix $q=(X,\theta)$ with $X=(x,y)\ne0$, $0<|\theta|<\pi$ and
$x\cos(\theta/2)+y\sin(\theta/2)\ne0$.  Let $I$ be a maximal
admissible component of its scalar inverse.  For every $\ell>0$,
$\{a\in I:N(a)\ge\ell\}$ is a compact subset of $I$, possibly empty.
\end{lemma}
\begin{proof}
Use the principal axis angle $a$, with $b=a-\theta$, on
$J=(-\pi/2,\pi/2)\cap(\theta-\pi/2,\theta+\pi/2)$, and write
$A_q=x\sin a-y\cos a$, $B_q=x\cos a+y\sin a$.
The circle equations give
\[
 m=\sin^2a+c_0^2=\sin^2b+c_1^2,\qquad
 \sin\theta\sin(2a-\theta)=-B_q(c_0+c_1).
\]
At least one of $|a|,|b|$ is at least $|\theta|/2$, so
$m\ge m_h:=\sin^2(\theta/2)>0$.  The defining integral makes $L$
strictly increasing, with limits $0$ and $+\infty$ at $0$ and $1$
\cite[Sections 19.2, 19.12]{DLMF}.
On a nonempty positive superlevel,
\[
 \ell L(m)\le A_q-Q_0<A_q\le R,\qquad
 m_h\le m<M_\ell:=L^{-1}(R/\ell)<1,\qquad
 A_q>\ell L(m_h)>0.
\]
Take a sequence in this superlevel converging in $\overline J$.
A limit at an end of $J$ would force $m\ge1$, since $a$ or $b$
would equal $\pm\pi/2$.  At an interior limit, $A_q$ stays positive.
If $B_q\to0$, boundedness of $c_i$ and the second circle identity
would give $\sin(2a-\theta)=0$.  On $\overline J$ one has
$-\pi<2a-\theta<\pi$, hence $a=\theta/2$, contradicting the
target's off-seam hypothesis.  Thus $B_q\ne0$ at the limit, and
the algebraic inverse extends there with $0<m<1$.
Local phase lifts make $Q_0$ continuous: the two endpoint phases
cannot coincide modulo $4K$, since their principal angles differ
by $\theta\ne0$.  Hence the limit is admissible with $N\ge\ell$.
A component is closed relative to the admissible domain, so the
limit belongs to $I$.  This proves compactness; the parameters
$z=\tan a$ and $\nu$ give the same conclusion.
\end{proof}

\begin{corollary}[Counts at positive levels]\label{cor:verified-counts}
Let $q$ and its maximal admissible component $I$ satisfy
Lemma~\ref{lem:positive-superlevels}.  Assume that every critical
point of $N$ with $N\ge\tfrac12$ and every
interior local minimum of $N$ satisfies \eqref{eq:verified-predicate}.
Off $\mathcal S\cup\mathcal C_1\cup\mathcal C_2\cup\mathcal C_3\cup\mathcal C_4$
this is the certified statement; on those sets it is an assumption.  Then
every critical point of $N$ on $I$ at a level $\ge\tfrac12$ is a strict
maximum, at most one critical point of $N$ on $I$ lies at a level $\ge1$,
and for every real $\ell\ge1$ the superlevel set
$\{\nu\in I:N(\nu)\ge\ell\}$ is an interval (possibly empty).
Writing $M=\sup_I N$, for every real $\ell\ge1$ the number of
distinct solutions of $N=\ell$ is $2$, $1$ or $0$ according as
$\ell<M$, $\ell=M$ or $\ell>M$.  In particular the count for
integer $n\ge1$ is at most two and non-increasing in $n$.
\end{corollary}
\begin{proof}
By \eqref{eq:verified-predicate} a critical point with $N\ge\tfrac12$ has
$N''<0$ and is a strict maximum, and an interior local minimum has
$N<\tfrac12$ and $R<L(m)$.  Let $a<b$ lie in $I$ with
$N(a),N(b)\ge\ell\ge1$.  If $N$ dips below $\ell$ between them,
its minimum on $[a,b]$ occurs at an interior point $d$.  By
Lemma~\ref{lem:modulus}, $m(d)\le\max(m(a),m(b))=:m_{\max}$.  At either
endpoint $j\in\{a,b\}$,
$L(m(j))\le N(j)L(m(j))=A_q(j)-Q_0(j)\le R$.  Consequently
$R\ge L(m_{\max})\ge L(m(d))$, contradicting the predicate at $d$.
Every superlevel set at a level $\ell\ge1$ is therefore an interval,
possibly empty.  Two strict maxima with values at least $1$ would
enclose an interior local minimum; the same inequality, using those
maxima as $a,b$, excludes it.  Hence there is at most one such
critical point.  If $N=\ell\ge1$ had three distinct solutions
$a<b<c$, superlevel connectedness would give $N\ge\ell$ on $[a,c]$,
making $b$ a local minimum, a contradiction.  This proves the
two-solution bound without assuming the existence of a maximum.
Now $M\le R/L(\sin^2(\theta/2))$ is finite.  If $M>0$, a maximizing sequence
eventually lies in the compact superlevel at $M/2$, so $M$ is
attained.  For $1\le\ell<M$, compactness of the superlevel and
continuity give a crossing on each side of the maximum; the
two-solution bound excludes any more.  At $\ell=M\ge1$ there is
exactly one maximizer, by the critical-point conclusion, and above
$M$ there are no roots.  If $M<1$ there are no positive integer
roots.  This proves the stated count law.
The inequality at $a,b$ needs $N\ge1$; it does not exclude two
maxima whose values lie in $[\tfrac12,1)$.
\end{proof}

\begin{remark}
The maximal-component and target off-seam hypotheses are essential:
a truncated interval or a chart component of a seam target can have
a positive level at an excluded end.  The monotonicity conclusion
starts at $n=1$.  It does not extend from $n=0$ to $n=1$: in the
far-field regime of \cite[Theorem 3.3]{SE2Counts}, one component
has one level-zero root and two roots at level one once its maximum
exceeds one.  The level-zero clauses therefore remain separate.
The corollary requires the endpoint-data hypothesis at all of its
specified critical points; membership of a target in a spatial region
alone does not establish that hypothesis.  Some necessary constraints are
explicit: along a component $m\ge\sin^2(\theta/2)$, so $m\ge0.02$ when
$|\theta|\ge0.284$; and a critical point with $m>0.98$ and
$N\ge\tfrac12$ needs $R>L(0.98)/2$, since
$NL=A_q-Q_0\le R$.  On $m\le0.98$, an amplitude increment of size
$\epsilon$ gives an elliptic phase increment between $\epsilon$ and
$\epsilon/\sqrt{1-0.98}$, by the integral defining $F$.  Thus
$\Delta\varphi$ within $0.02$ of $0$ or $2\pi$ puts the phase increment
within $0.02/\sqrt{0.02}<0.142$ of zero or one full period, respectively.
The slab $\mathcal S$ is a neighbourhood of the seam $B_q=0$, off which
Proposition~\ref{prop:level} is stated.  The corner and tube samples do
not prove exclusions on those continuous sets.  On the finite family the
reported largest values at an interior minimum are approximately $0.45$
for $R/L(m)$ and $0.39$ for $N$ on the finest grid of
\path{research/e78_statement_L.py}, and up to $0.43$ for both on a
coarser grid, so the margin to the threshold $\tfrac12$ of
\eqref{eq:verified-predicate} is about
$0.05$; these are numerical observations, not uniform bounds.  Proposition~\ref{prop:far-incidence} separately
covers fixed far-field rays.  Level $0$ is outside the positive-level
corollary.
\end{remark}

\paragraph{Method.}  Every quantity of Proposition~\ref{prop:N-prime} is
evaluated in ball arithmetic \cite{Johansson2017Arb,Tucker2011}
(\texttt{arb}, through \texttt{python-flint}) as a first-order Taylor
model over a box: centre value, box enclosure, and box enclosure of the
gradient with respect to $(m,\mu,\Delta\varphi)$, the enclosure being
tightened by the mean-value form after every operation.  The incomplete
integrals are enclosed by monotonicity: $F$ and $E$ increase in
$\varphi$, and in $m$ the derivatives
$\partial_mF=\tfrac12\int_0^\varphi\sin^2\theta\,(1-m\sin^2\theta)^{-3/2}d\theta$
and
$\partial_mE=-\tfrac12\int_0^\varphi\sin^2\theta\,(1-m\sin^2\theta)^{-1/2}d\theta$
have the signs of $\varphi$ and of $-\varphi$, so the extremes of each
over a box are corner values, and the union of the four corner
evaluations is the enclosure for either sign of $\varphi$ (the lower
endpoint phase $\mu-\Delta\varphi/2$ is negative on most of
$\mathcal B$); the derivatives are in closed form.  Two reformulations make the evaluation regular: with
$\Gamma=c_1s_0d_0-c_0s_1d_1$, the polynomial
$\Lambda=(\Gamma^2-m(1-m)B_q^2)/(c_0c_1)$ and $\mathcal W=\Delta E+\beta_N(A_q-Q_0)$,
the critical condition is the quadratic
\[
 \Psi(A_q)=-c_0c_1\beta_NA_q^2+bA_q+c=0,\qquad
 b=-\Gamma-c_0c_1(\Delta E-\beta_NQ_0)+\beta_N\Gamma,
\]
\[
 c=\Lambda+\Gamma(\Delta E-\beta_NQ_0),
\]
in $A_q$, with $\Psi=2B_q(1-m)L(m)\,\partial_\nu N=\Lambda-A_q\Gamma-(c_0c_1A_q-\Gamma)\mathcal W$,
so that at a root $\operatorname{sgn}N''=\operatorname{sgn}(B_q\,\partial_\nu\Psi)$.
For the root test a polynomial form avoids division at an endpoint flip.
Put $\mathcal F=c_0c_1A_q-\Gamma=B_q\,\partial_\nu m$ and
$\mathcal W=\Delta E+4E(m)N$.  At $\Psi=0$, hence $\partial_\nu N=0$,
\[
 B_q\,\partial_\nu\mathcal W
 =\frac{d_1R_1-d_0R_0-\mathcal F A_q}{2m},
 \qquad \partial_\nu\varphi_i=\frac{R_i}{2mB_q},
\]
where the $R_i$ are the polynomial amplitude numerators of
\path{research/e78_regular_prep.py}.  Substituting this expression when
differentiating $\Psi$ gives a cubic $P_{\rm crit}(A_q)=p_3A_q^3+p_2A_q^2+p_1A_q+p_0$ with polynomial
coefficients in $(s_i,c_i,d_i,\Delta E,\beta_N,Q_0)$ and
\[
 2mB_q\,\partial_\nu\Psi
 =P_{\rm crit}(A_q)-\frac{m\beta_N\mathcal F}{1-m}\Psi(A_q).
\]
This identity follows also from the omitted term
$4E(m)\partial_\nu N$ in $\partial_\nu\mathcal W$ and is checked exactly
modulo the endpoint circle identities.  Thus $P_{\rm crit}$ equals
$2mB_q\,\partial_\nu\Psi$ at every critical root, including flips by
continuity, but is not the unrestricted derivative away from those roots.
Its leading coefficient is
\[
 p_3=c_0^2c_1^2-m\beta_N(c_0^2+c_1^2).
\]
The unrestricted cubic in \path{research/e78_cubic_prep.py} includes
$2m\beta_N'(m)=-1-m\beta_N^2/(1-m)$ and has the additional leading term
$m\beta_N^2c_0^2c_1^2/(1-m)$.  The earlier polynomial is valid for the
critical-root test; interpreting it as the unrestricted derivative was
incorrect.  The verifier uses the polynomial coefficients generated by
\path{research/e80_critical_cubic_prep.py}.  When
$p_3<0$ on a box, the \emph{cubic threshold}
\[
 A_{\rm th}=1+\frac{|p_2|+|p_1|+|p_0|}{-p_3},
\]
formed from the bounds of the enclosures, gives $B_q\,\partial_\nu\Psi<0$,
hence $N''<0$, at every root $A_q\ge A_{\rm th}$: for $A_q\ge1$,
the sum of the lower-order terms is at most
$(|p_2|+|p_1|+|p_0|)A_q^2$, while
$(-p_3)A_q>|p_2|+|p_1|+|p_0|$.  When the enclosure and threshold comparisons succeed, this settles the
root that escapes to infinity at a flip; otherwise the box remains undecided.  Let
$\mathcal E=\{N<\tfrac12,\ R<L(m)\}$ be the \emph{escape set}, the
second alternative of \eqref{eq:verified-predicate}.  A box passes when
its enclosures establish that $\Psi$ has no real root, or that every root
$A_q$ of $\Psi$ on the box satisfies
\begin{equation}\label{eq:box-predicate}
 \begin{gathered}
 A_q\le0,\qquad\text{or}\qquad B_q\,\partial_\nu\Psi<0,\qquad\text{or}\qquad
 A_q\ge A_{\rm th},\\
 \text{or}\qquad N<\tfrac12\ \text{ and }\ R<L(m)\ \text{ at the root}
 \end{gathered}
\end{equation}
(a root that escapes to infinity because the leading coefficient of
$\Psi$ straddles zero passes when its modulus is at least $A_{\rm th}$; when the
enclosure of the discriminant of $\Psi$ straddles zero, its square root
has no first-order model and is taken as the plain interval
$[0,\sqrt{\max\operatorname{disc}}]$, the roots and every quantity
computed from them then being carried as plain intervals without
mean-value tightening); the
first alternative is an inadmissible root, the second and third give
$N''<0$, the fourth is the escape set, so a passed box satisfies
\eqref{eq:verified-predicate}.  Otherwise the box is bisected along its
widest side, down to width $\min(2\cdot10^{-3},m/25)$ in the first pass
and $2.5\cdot10^{-4}$ in the refinement.  The removal of the slab
$\mathcal S$ and of $\mathcal C_1,\mathcal C_2,\mathcal C_3$ is implemented
by passing without evaluation every box that lies inside one of them
(the verifier records them as deferred); these sets are outside the
scope of the proposed certificate.  The fifth set $\mathcal C_4$ is
declared after the run as the union of the cells whose budget passes left
boxes open; its cells were evaluated like all others, their accepted
leaves are kept, and no box outside $\mathcal C_4$ remains open.
The certificate run covers the initial grid of $24\times20\times80$
cells, one checkpoint per cell, with four budget passes of $256$, $4{,}096$,
$16{,}384$ and $65{,}536$ evaluations; a box is accepted by the first of
four reviewed predicate methods (the scaled critical-root test above and
three auxiliary root-branch tests; \path{research/e80_combined_predicate.py}),
bisected along its widest side otherwise, and left open when
the budget of its cell is exhausted.  Boxes inside the four deferred sets
are recorded as deferred, not counted as verified, and are outside the
scope of the certificate; no box reached the minimum width.  The
$38{,}400$ checkpoints carry the source hashes and the canonical
subdivision trace of their cell; an independent semantic replay re-proves
every accepted leaf, re-walks every partition and checks each resumed
pass against its parent, with no failure: $1{,}072{,}966$ accepted leaves,
$15{,}777$ deferred leaves, $304$ open leaves (all in $\mathcal C_4$)
(\path{data/e80_full_grid_pass1.json}, \path{data/e80_full_grid_audit.json},
\path{data/e80_full_grid_open_cells.json}).  The old run statistics and
residual clusters of the defective verifier are withdrawn as certificate
evidence and retained only in the research notebook; volume fractions use
the ambient rectangle before exclusions as denominator.
Scripts:
\path{research/e78_verify5.py}, \path{research/e78_tm.py},
\path{research/e78_regular_prep.py}, \path{research/e78_cubic_prep.py},
\path{research/e80_combined_predicate.py}, \path{research/e80_combined_checkpoint.py},
\path{research/e80_full_grid_local.py}, \path{research/e80_full_grid_audit.py}.

\begin{conjecture}[Two maxima]\label{conj:two-maxima}
For every target off the seams with two admissible intervals and every
integer $n\ge0$, the equation $N=n$ has at most two solutions on each
interval; for $n\ge1$ their number is non-increasing in $n$, and at levels $\ge1$
the level function has on each interval at most one critical point, a
maximum; at level $0$ no critical point other than the maximum of the
interval lies on the level, so that the conjugate preimages of a target
are its interval maxima at integer levels; the two maxima have different
residual types, $\delta>2K$ at one and $\delta<2K$ at the other; and the
maximum of right type exceeds the maximum of left type, $M_R>M_L$.
Corollary~\ref{cor:verified-counts} proves the bound of two solutions,
uniqueness of a critical point, and count monotonicity for $n\ge1$
on maximal off-seam components under its explicit sign-predicate
assumption.  The repaired verifier has not yet established that
assumption on a certified domain.  The unconditional positive-level
claims, the level $0$ (its count and critical-point clause),
the removed sets and margins of $\mathcal B$,
and the two clauses on the types remain conjectural.
\end{conjecture}

\paragraph{Evidence.}  A replay with $1{,}000$ accepted random targets
in $0.5<|X|<14$, $0.05<|\theta|<\pi-0.05$ and at least one finite
admissible interval gives $1{,}798$ finite intervals: the sampled detector
reports one nonnegative-level maximum on $1{,}667$ and none on $131$.
A separate sample of $600$ targets with two finite intervals gives $594$
with both interior maxima detected; their types differ, and $M_R-M_L$
ranges from $0.1396$ to $274.913$, below $1$ on $94$ targets
(\path{research/e76_unimodality_scan.py}, seed $20260907$, argument $1000$;
\path{research/e76_two_maxima.py}, seed $7$, argument $600$).
These finite-grid results do not exclude missed critical points.

A historical $60\times120\times120$ endpoint-data scan reported
$405{,}558$ quadratic roots with $A_q>0$, $N\ge0$: $402{,}796$ maxima
and $2{,}762$ minima.  A portable replay using a separately derived
critical-curvature formula reproduces these finite-grid counts with
both the original and a cancellation-resistant quadratic formula
(\path{research/e80_e77_replay.py}); all $2{,}762$ candidate targets
match the saved survey in order, with spatial-radius error below
$1.9\times10^{-12}$.  Selected targets also pass independent
$70$-digit differentiation checks.  The original scratch polynomial
remains unavailable; this replacement does not give an interval
certificate or certify every extremum detected by the subsequent
chart scans.  The saved survey contains $2{,}762$ records marked as
minima by those scans.  Their $|\theta|$ is at most
$0.1597$, $|X|$ at most $1.482$, and minimum level at most $0.2743$;
the smaller recorded neighboring maximum is at most $0.2744$.
No record flags four crossings of an integer level.  These saved chart observations do not certify every detected
stationary point or exclude additional unsampled extrema.  The
finite-difference detector is unreliable for the shallowest dips:
on six reconstructed targets with $5$--$11$ historically reported
critical points, $30$ of $81$ examined sampled derivative signs are
incorrect in binary64, while $70$-digit recomputation of the same
stencils removes those mismatches.  Local interval checks isolate
two maxima and one minimum for each of these six targets, all at
levels strictly between $0$ and $1$; they do not establish completeness
outside the $18$ isolating boxes
(\path{references/review_2026-09-21_neighbors/}).  The finer grid discussed after
Corollary~\ref{cor:verified-counts} reports levels up to $0.39$.
A separate saved small-heading survey records $86$ dips, none flagged
as crossing an integer level four times
(\path{research/e77_critical_quadratic.py},
\path{research/e77_dips_survey.py}, \path{research/e77_dips_neighbors.py},
\path{data/e77_dips_survey.json}).  The locally certified positive-level minima
rule out all-real-level unimodality; the integer
statement remains conjectural.

\begin{corollary}[Conditional on Conjecture~\ref{conj:two-maxima}]\label{cor:incidence}
Let $q=(X,\theta)$ satisfy $X\ne0$, $0<|\theta|<\pi$, be off the
scalar-inverse seams, and have two admissible intervals as in the
conjecture.  Such a target has at most two rotating conjugate preimages.
Within this target domain the left-sheet images are pairwise disjoint,
and so are the right-sheet images.  For the octant-matching $g$, with
the intersections restricted to this domain,
\[
 F(n,-)\cap g\,F(j,+)=\{q:\ M_L(q)=n-1,\ M_R(q)=j\},
\]
which is empty unless $j\ge n$.  The number of rotating geodesics through
a target is determined by $\lfloor M_L\rfloor$, $\lfloor M_R\rfloor$ and the
end values of $N$ on the two intervals; it is constant on the chambers
cut by the sheets and changes by two across a sheet.
\end{corollary}
\begin{proof}
By Proposition~\ref{prop:level} conjugate preimages are critical points
at integer levels; with one maximum per interval at levels $\ge1$ and
no critical point at level $0$ other than the interval maxima, there
are at most two, of different types, and a target on $\Sigma_n^-\cap\Sigma_{n'}^-$ would need
$M_L=n-1=n'-1$.  The incidence set is the double integer condition, and
$M_R>M_L$ forces $j>n-1$.  With at most two solutions per integer level,
non-increasing at positive levels, the number of integer crossings is read off
from the maximum and the end values, and crossing a sheet changes the
count by two (a maximum passing an integer level creates or destroys a
pair).
\end{proof}

In the far field the unimodality part is a theorem at every real level.

\begin{proposition}[Far-field incidences]\label{prop:far-incidence}
Fix a spatial direction $a_0$ and a heading $\theta\ne0$ with
$b_0=a_0-\theta$, both in $(-\pi/2,\pi/2)$ and with
$\sin^2b_0\ne\sin^2a_0$, and let $q_R=(R\,e(a_0),\theta)$ with
$e(a_0)=(\sin a_0,-\cos a_0)$ and $L_*=L(\max(\sin^2a_0,\sin^2b_0))$ as in
\cite[Section~3]{SE2Counts}.  There is
$R_0$ such that for $R\ge R_0$: the level function of $q_R$ has exactly
one critical point on each of its two admissible components, a
nondegenerate maximum, of heights $M_\pm(R)=R/L_*+O(1)$; hence $q_R$
has at most two conjugate preimages, it lies on the conjugate locus
exactly when $M_-(R)$ or $M_+(R)$ is an integer, on two sheets exactly
when both are, and the two sheets through such a point belong to cells
whose indices differ by $O(1)$ uniformly in $R$.
\end{proposition}
\begin{proof}
The first assertion is \cite[Theorem 3.3]{SE2Counts}, and the exclusion
of the sign $s_2=-1$ at large $R$ is \cite[Section~3]{SE2Counts}; the rest is
Proposition~\ref{prop:level}, the level of a conjugate preimage being
$n-1$ or $n$ for a sheet of cell $n$.
\end{proof}

What is proved here is therefore the octant restriction, the
level-function characterization, the conditional count implication at
levels $\ge1$ (Corollary~\ref{cor:verified-counts}), and the far-field
structure along fixed rays; the map for the first four cells is computed
numerically only; the global structure at all radii, and the inequality
$M_R>M_L$ behind the rule $j\ge n$, rest on
Conjecture~\ref{conj:two-maxima}.

\section{Numerical verification}\label{sec:numerics}

The accompanying numerical scripts are in the
\href{https://github.com/moiseevigor/se2-conjugate-locus}{project repository}, folder
\path{research/}, prefixes \texttt{e74}--\texttt{e78}.
The historical summaries below describe finite-sample diagnostics;
the complete collection of these runs has not been independently
replayed.  At $84$ points of the open sheet, six moduli including
$0.9089$ (near, but not equal to, $k_0$) and $0.99$, both branches,
the independently replayed finite-difference
differential of \eqref{eq:exp} at $40$ digits has singular values
$\sigma_2/\sigma_1\ge4.9\times10^{-3}$, $\sigma_3/\sigma_1\le6\times10^{-28}$,
kernel time component $\le7\times10^{-25}$, and the kernel derivative of
$\det d\Exp$ has the sign predicted by \eqref{eq:fold-identity} at every
point.  The replayed minimum of $\sigma_2/\sigma_1$ is
$0.004990276288\ldots$, correcting the previously rounded lower bound.
Independent $60$-digit differentiation at eight selected points, including
both phases near this minimum, agrees with separately integrated
Hamiltonian and variational equations: the largest endpoint error is
below $9\times10^{-14}$, and the largest differential error, normalized
by $1+\max|D\Exp|$, is below $4\times10^{-13}$
(\path{research/e80_geometry_audit.py},
\path{references/review_2026-09-21_geometry_diagnostics/}).
These are finite-sample diagnostics, not interval enclosures.
The remaining historical summaries have not all been independently
replayed.  The symbolic $\theta_k$, $\theta_\psi$, $v(J_2)$ agree with finite
differences to $3\times10^{-20}$.  At the Maxwell edge and at $(2K,K)$, for
four and five moduli, $|v(J_2)|\le3\times10^{-28}$, away from $k_0$ the second kernel derivatives agree
with Theorem~\ref{thm:edges} to ten and twelve digits and the rank condition
holds; at $k_0$ the singular values are $1$, $8\times10^{-22}$, $0$.
On the $\alpha_1$ edge $Q\in[9\times10^{-6},8\times10^{-3}]$ and $|\mathsf Z|\le3\times10^{-42}$
at four moduli.  The atlas chart agrees with the raw map to $2\times10^{-28}$ at
$250$ points; $H_p\ge0.0239$, $a'\ge0.0465$, $\zeta'\ge0.003$ there; $H_1$
increases on $92$ moduli from $0.02$ to $19.6$; $U$ increases along the
slope levels $H=0.05,0.3,1,3,10$; the lower seam is ordered; and the
escape bound holds on the strips $k=0.9,0.99,0.999,0.999999$, where the
minimal spatial radius rises from $3.26$ to $14.18$.
Independent replay covers the edge and chart samples.  A separate
$50$-digit implementation checks the upper slopes, lower seams, five
nonempty slope levels and four escape strips, with selected $60$-digit
raw-map derivative checks
(\path{references/review_2026-09-21_edge_embedding/}).  All bounds remain
finite-sample diagnostics.  On the later sheets,
at $180$ points of cells $1$--$5$, both brackets and both branches, three
moduli including $0.9089$, the identity \eqref{eq:fold-identity} holds to
$10^{-8}$ and the corank is one; $nQ$ divided by the leading term of
\eqref{eq:Q-leading} tends to one with an $O(1/n)$ deviation up to
$n=80$; the seam script tests $30$ pairs $(k,n)$ and prints ten of them,
reporting agreement with \eqref{eq:beta-seam} to $10^{-18}$ and
$d(kb_n)/dk>0$.  The derivative $\partial_k^{\rm src}\beta_1$ instead
has sign $(-1)^n$, by \eqref{eq:beta-seam}; independent $60$-digit checks
at $k=0.6,0.95$ and $n=1,2,3$ confirm this sign and the formula.
The
intrinsic umbilic test gives $D_4^+$ at the central points of cells
$1$--$3$ (\path{research/e75_all_sheets.py}, \path{research/e75_step2_seams_all_n.py}).
These later-sheet scripts and the $48$ leading-term samples were replayed;
a separate $60$-digit implementation checks all $180$ open-sheet points,
critical Hessians at cells $1,2,3,4,6$ and six gap examples
(\path{references/review_2026-09-21_later_sheets/}).

The independent September 2026 review checked the exact fold certificates,
atlas identities and umbilic jets, replayed the saved incidence targets
and the two random-target summaries above, and checked selected raw
endpoints and the scalar-inverse dependency at $60$ digits.  Its logs,
source hashes and limitations are retained in
\path{references/review_2026-09-20_caustic/} and
\path{references/review_2026-09-21_caustic_provenance/}.
The interval sign predicate still requires a complete corrected replay
and an independent coverage and residual audit.

\section{Relation to prior work}

Sachkov \cite{Sachkov2010Conjugate} proves that the first conjugate time
lies in $[2kp_0,4kK]$, states the bracket $[2kp_0,2kp_1]$ and gives the
Jacobian \eqref{eq:j2};
Moiseev--Sachkov \cite{MoiseevSachkov2010Maxwell} give the exponential map
\eqref{eq:exp} and the half-angle chart; \cite{SE2FirstConjugate} makes the
first conjugate time exact and supplies the Riccati identities
\eqref{eq:riccati}.  The Lagrangian classification, with the
determinacy and versality of the generating families, is Arnold's
\cite{Arnold1972NormalForms}; the fold and cusp criteria are Morin's
\cite[Chapter~VII]{GolubitskyGuillemin1973}; Itoh--Kiyohara
\cite{ItohKiyohara2021} describe the conjugate locus of an ellipsoid with
the same three local types.  The
original derivations of the atlas, the univalence argument and the
closed-incidence argument were made in the half-angle chart; the present
paper re-derives the corank, kernel, fold identity and edge classification
in raw coordinates by the first variation, which shortens them and makes
the fold factor explicit.  The local structure of every later sheet is
Section~\ref{sec:every-sheet}, whose certificates and seam inequalities
are reproduced in Appendix~\ref{app:proofs}, and the incidences between
sheets are Section~\ref{sec:incidence}; the
multiplicity of every point follows from Conjecture~\ref{conj:two-maxima}
and is the subject of \cite{SE2Counts}.

\section*{Code and data availability}

All scripts, the certificate tables and the figure code are in the
repository cited above: \path{research/e74_p8_kernel_rederivation.py},
\path{research/e74_step2_fold_identity.py}, \path{research/e74_ring.py},
\path{research/e74_step3_edges.py}, \path{research/e74_step4_d4.py},
\path{research/e74_step5_embedding_audit.py}, \path{scripts/caustic_figures.py},
\path{scripts/review/verify_q_certificate_fractions.py},
\path{scripts/first_caustic_projection_reduction.py},
\path{scripts/universal_right_fold_proof.py} (run with
\texttt{PYTHONPATH=.:src}, as it imports the \path{scripts} package),
\path{research/e78_verify5.py}, \path{research/e78_tm.py},
\path{research/e78_regular_prep.py}, \path{research/e78_cubic_prep.py},
\path{research/e78_statement_L.py}, \path{research/e75_leading_term.py},
\path{research/e75_step3_certificate_domains.py},
\path{scripts/first_caustic_global_univalence_proof.py},
\path{research/e79_atlas_identity.py}.  The certificate table
\path{data/open_sheet_q_certificate.json} records the SHA-256 digest
\texttt{7e9fabf7\ldots950c} of its canonical coefficient table in its
\texttt{sha256} field, and \path{data/universal_right_sv_certificates.json}
likewise \texttt{3643ff00\ldots13ff07} (digests of the tables, not of the
files).  The prior interval summaries
\path{data/e78_verify_L.json} and \path{data/e78_refine_L.json}
are historical outputs of the defective verifier, not proof evidence.
Independent primitive regression checks are recorded in
\path{data/e80_primitives_check.json}; the regenerated certificate and its
replay are \path{data/e80_full_grid_pass1.json},
\path{data/e80_full_grid_audit.json} and
\path{data/e80_full_grid_open_cells.json}.  The long
working manuscript \path{manuscripts/first_conjugate_time_c2.tex} of the
repository (the research compendium) contains the original derivations in
the half-angle chart; the three proofs of Appendix~\ref{app:proofs} are
transplanted from it.

\section*{AI assistance}

This paper was prepared with the assistance of Claude (Anthropic), a large
language model used through the Claude Code agent in September 2026.  The
model re-derived the kernel, the fold identity and the edge classification
in raw coordinates from the author's working manuscript, found the factor
in \eqref{eq:fold-identity} numerically before proving it, wrote the
verification and figure code, and drafted the text.  Every statement was
checked by the author, who takes full responsibility for the content; no
claim rests on the model's output alone.  OpenAI Codex assisted with the
September 2026 adversarial review, independent certificate and endpoint
checks, and the critical-root polynomial reformulation.

\appendix
\section{Three proofs from the working manuscript}\label{app:proofs}

The three results below are used in Proposition~\ref{prop:atlas},
Lemma~\ref{lem:upper-slope} and Theorem~\ref{thm:every-sheet}.  They are
transplanted, in the notation of this paper, from the working manuscript
named in the code and data section.

\begin{proposition}[Reduction of the projection determinant]\label{prop:app-projection}
In the amplitude chart $(A,B,k)=(\am p,\am\tau,k)$ of the source put
\[
 \Delta=\det\frac{\partial(U,V,J_2)}{\partial(A,B,k)} .
\]
On the fundamental patch $0<B<\pi/2$, $(J_2)_B=2\sin B\cos B\,(\alpha-\beta)>0$,
the restriction of $(U,V)$ to the surface $J_2=0$ has
\begin{equation}\label{eq:app-projection-restriction}
 \det\frac{\partial(U,V)}{\partial(A,k)}\Big|_{J_2=0}=-\frac{\Delta}{(J_2)_B},
\end{equation}
and on $J_2=0$, in the variables $g=-c$, $\mathsf x=\mathsf u/g$,
$\mathsf t=L_0\mathsf v/\mathsf u$, $\mathsf y=g^2$ of Theorem~\ref{thm:fold},
\begin{equation}\label{eq:app-projection-reduction}
 \Delta=\Pi\cdot g^3\mathsf x^2\,
 \bigl[L_0^2+L_0\mathsf y(1-\mathsf t+\mathsf t\mathsf x)+L_0\mathsf t+\mathsf y\mathsf t\mathsf x\bigr]\,
 P(L_0,\mathsf x,\mathsf t,\mathsf y),
\end{equation}
with $P$ the certificate polynomial of \eqref{eq:certificate} and $\Pi$ a
product of factors of fixed sign on the patch.  Consequently
$\partial(U,V)/\partial(z,k)\ne0$ on every open patch, and the sign in
\eqref{eq:orientation} is constant on the fundamental patch.
\end{proposition}
\begin{proof}
Expanding $\Delta$ along its third row after the column operations that
kill $(J_2)_A$ and $(J_2)_k$ gives \eqref{eq:app-projection-restriction},
because on $J_2=0$ one has $B_A=-(J_2)_A/(J_2)_B$ and $B_k=-(J_2)_k/(J_2)_B$.
The reduction \eqref{eq:app-projection-reduction} is an exact computation
in the ring of Section~\ref{sec:fold}: $\Delta$ is formed from
\eqref{eq:atlas-target}, \eqref{eq:j2} and the modulus derivatives of
Theorem~\ref{thm:fold-identity}; the factor $C=\cos B>0$ is divided out;
$J_2=0$ is imposed through $S^2=r$; the Riccati variables enter by
$q=-ds/(a+c)$, $e=ds/(b+c)$ (that is, \eqref{eq:a-below-g} and its
analogue for $b$); the result is reduced modulo $c^2+s^2=1$,
$d^2+ms^2=1$; the fixed factor $a\,d\,k^2s^7(k^2-1)(ks-1)^4(ks+1)^4$ is
removed; and the substitution $c=-g$, $a=g(1-\mathsf x)$,
$b=g(1+\mathsf x\mathsf t/L_0)$, $m=1/(1+L_0)$, $s^2=1-g^2$ leaves
$g^3\mathsf x^2$ times the bracket of \eqref{eq:app-projection-reduction}
times a polynomial in $(L_0,\mathsf x,\mathsf t,g^2)$, which is verified to
coincide with $P$ term by term.  The script
\path{scripts/first_caustic_projection_reduction.py} carries out this
chain in exact arithmetic and raises an error if the final polynomial
differs from $P$.  The bracket is positive because $0<\mathsf t<1$, and
$P>0$ by Theorem~\ref{thm:fold}; so $\Delta\ne0$, hence
$\partial(U,V)/\partial(A,k)\ne0$ on $J_2=0$ by
\eqref{eq:app-projection-restriction}, and $\partial(U,V)/\partial(z,k)$
differs from it by the factor $\partial A/\partial z=\partial_z\am(\mathcal P(z,k),k)>0$,
since $\mathcal P$ increases with $z$ by the first-time theorem.  The other
patches follow by the sign actions of Proposition~\ref{prop:atlas}.
\end{proof}

\begin{lemma}[Three bounds at the supercritical endpoint]\label{lem:app-bounds}
Let $k_0<k<1$, $m=k^2$, $w(\theta)=\sqrt{1-m\sin^2\theta}$, and let
$\varpi=2K-p_\alpha\in(0,K)$, $\tilde e=2E-\epsi(\varpi)$ be as in the
proof of Lemma~\ref{lem:upper-slope}.  Then
\[
 \text{(i)}\ k_0^2>\tfrac34,\qquad
 \text{(ii)}\ \tilde e>\tfrac32,\qquad
 \text{(iii)}\ 2(1-m)(K-E)<\tfrac34 .
\]
\end{lemma}
\begin{proof}
(i) Pair $\theta$ with $\pi/2-\theta$ in the complete integrals.  From
$w(\theta)^2w(\pi/2-\theta)^2=(1-m\sin^2\theta)(1-m\cos^2\theta)
=1-m+m^2\sin^2\theta\cos^2\theta\ge1-m$,
\[
 K=\int_0^{\pi/4}\frac{w(\theta)+w(\pi/2-\theta)}{w(\theta)\,w(\pi/2-\theta)}\,d\theta
 <\frac1{\sqrt{1-m}}\int_0^{\pi/4}\bigl[w(\theta)+w(\pi/2-\theta)\bigr]d\theta
 =\frac{E}{\sqrt{1-m}}
\]
for every $0<m<1$.  At $m=3/4$ this is $2E>K$.  Since $2E-K$ decreases
strictly in $m$ ($E$ decreases and $K$ increases) and vanishes exactly at
$m=k_0^2$, we get $k_0^2>3/4$.

(ii) With $\tilde s=\sn\varpi$, $\tilde g=\cn\varpi$, $\tilde d=\dn\varpi$,
the endpoint equation $\alpha_1(p_\alpha)=0$ reads, as in the proof of
Lemma~\ref{lem:upper-slope},
$q(\tilde e\tilde s-\tilde g\tilde d)+\tilde d(\tilde g\tilde e+\tilde s\tilde d)=0$
with $q,\tilde d,\tilde g,\tilde e,\tilde s>0$, so $D_0(\varpi)<0$ for
$D_0(x)=e(x)\sn x-\cn x\dn x$, $e(x)=2E-\epsi(x)$.  Differentiating with
$e'=-\dn^2$ gives $D_0'(x)=\cn x\,\bigl[e(x)\dn x+m\sn x\cn x\bigr]>0$ on
$(0,K)$.  Let $\phi_0=\pi/6$ and $x_0=F(\phi_0,k)$, so that
$D_0(x_0)=\tfrac12\bigl[2E-E(\phi_0\,|\,m)-\sqrt3\,w(\phi_0)\bigr]$.  Put
$h(\theta)=w(\theta)-\cos\theta$; then $h(0)=0$ and
$h'(\theta)=\sin\theta\,(w-m\cos\theta)/w>0$ because
$w^2-m^2\cos^2\theta=(1-m)(1+m\cos^2\theta)>0$, so $h$ increases and is
positive on $(0,\pi/2]$.  Subtracting the same expression at $m=1$, where
it equals $2-\sin\phi_0-\sqrt3\cos\phi_0=0$,
\[
 2E-E(\phi_0\,|\,m)-\sqrt3\,w(\phi_0)
 =\int_0^{\phi_0}h+2\int_{\phi_0}^{\pi/2}h-\sqrt3\,h(\phi_0)
 >\Bigl[2\Bigl(\frac\pi2-\phi_0\Bigr)-\sqrt3\Bigr]h(\phi_0)>0,
\]
since $2\pi/3>\sqrt3$.  Hence $D_0(\varpi)<0<D_0(x_0)$, so $\varpi<x_0$
and $\am\varpi<\phi_0$, and therefore
\[
 \tilde e=2E-E(\am\varpi\,|\,m)>2E-E(\phi_0\,|\,m)>\sqrt3\,w(\phi_0)
 =\sqrt3\sqrt{1-m/4}>\sqrt3\cdot\tfrac{\sqrt3}2=\tfrac32 .
\]

(iii) Since $1/w-w=m\sin^2\theta/w$,
$K-E=m\int_0^{\pi/2}\sin^2\theta\,w(\theta)^{-1}d\theta\le
m\,\pi/(4\sqrt{1-m})$, so $2(1-m)(K-E)\le\tfrac\pi2\,m\sqrt{1-m}$.  The
function $m\sqrt{1-m}$ decreases for $m>2/3$, and $m>k_0^2>3/4$ by (i),
so $m\sqrt{1-m}<\tfrac34\cdot\tfrac12$ and
$2(1-m)(K-E)<3\pi/16<3/4$.
\end{proof}

\begin{proposition}[$Q<0$ on the right sheets]\label{prop:app-right}
For every $n\ge1$ and $0<k<1$, $Q<0$ on the open sheet $\Sigma_n^+$.
For $k<k_0$ this includes its $\alpha_1$ end $ab=1$; the critical end
at $k=k_0$ has $Q=0$.
\end{proposition}
\begin{proof}
Every term of $Q$ has even total degree in $(a,b,c)$, so $Q$ is invariant
under the simultaneous negation of $(a,b,c)$ and we may work with the
sign-normalized variables of \eqref{eq:sheet-signs}: $a=g+\mathsf u$,
$b=g-\mathsf v$, $c=-g$, $g=|\cn p|\in(0,1)$, $\mathsf u=\rho/q$,
$\mathsf v=\rho/e$, $\rho=|\sn p|\dn p$, on $p=2nK+\xi$ with $0<\xi<K$.
On the sheet $b>0$, so $\mathsf w:=\mathsf v/g<1$; $0<ab\le1$ with
equality exactly at the $\alpha_1$ end; and $e=\epsi(p)>2nE>2$, so
$\mathsf v<\rho/2$.  Put $L_0=(1-m)/m$, $\mathsf x=\mathsf u/g$,
$\mathsf t=L_0\mathsf v/\mathsf u=(1-m)q/(me)$, $\mathsf T=2\mathsf t$ and
$\mathsf y=g^2$.

\emph{Step 1: $\mathsf t<1/2$ on every cell $n\ge1$.}  The inequality is
equivalent to $F(p)=mp-(2-m)q(p)>0$.  By quasi-periodicity
$q(p)=2n(K-E)+\xi-\epsi(\xi)$, so $F(p)=2nD+F_0(\xi)$ with
\[
 D=mK-(2-m)(K-E)=(2-m)E-2(1-m)K,\qquad
 F_0(\xi)=m\xi-(2-m)\bigl(\xi-\epsi(\xi)\bigr).
\]
Writing $D=m\int_0^{\pi/2}Y(\theta)/w(\theta)\,d\theta$ with
$Y=1-(2-m)\sin^2\theta$ and $w=\sqrt{1-m\sin^2\theta}$, and pairing
$\theta$ with $\pi/2-\theta$, where $Y$ becomes $m-Y$ and $w$ becomes
$B_0=\sqrt{1-m+m\sin^2\theta}$, the integrand of the pair is
$Y/A_0+(m-Y)/B_0$ with $A_0=w(\theta)$; on $0\le\theta\le\pi/4$ one has
$Y\ge m/2>0$, and when $Y>m$ the identity, with $x=\sin^2\theta\le1/2$,
\[
 Y^2B_0^2-(Y-m)^2A_0^2=m(2x-1)\bigl\{-(1-m)-x(1-x)(2-m)^2\bigr\}>0
\]
gives $Y/A_0>(Y-m)/B_0$.  So $D>0$.  On $[0,K]$, $F_0(0)=0$,
$F_0(K)=D>0$ and $F_0'(\xi)=m\{1-(2-m)\sn^2\xi\}$ changes sign once, from
positive to negative; hence $F_0\ge0$ and $F(p)=2nD+F_0(\xi)>0$.

\emph{Step 2: the margins.}  In these variables $\mathsf v<\rho/2$ and
$ab\le1$ read
\[
 \mathsf S_V=(1-\mathsf y)(L_0+\mathsf y)-4(1+L_0)\mathsf w^2\mathsf y>0,\qquad
 \mathsf S_Z=\mathsf T-\mathsf y(1-\mathsf w)(\mathsf T+2L_0\mathsf w)=2\mathsf t(1-ab)\ge0 .
\]
For fixed $(L_0,\mathsf x,\mathsf t)$ both cut upper intervals in
$\mathsf y$: $\mathsf S_V$ is concave in $\mathsf y$, equal to $L_0$ at
$\mathsf y=0$ and negative at $\mathsf y=1$, and $\mathsf S_Z$ is linear
decreasing in $\mathsf y$.

\emph{Step 3: concavity.}  Direct substitution gives
\[
 -Q=\frac{g^2\mathsf x}{L_0^3(1+L_0)}\,P^+(L_0,\mathsf x,\mathsf t,\mathsf y),
\]
where $P^+$ is a polynomial, quadratic in $\mathsf y$, with leading
coefficient $(L_0-\mathsf t\mathsf x)^2(1+\mathsf x)\{\mathsf t\mathsf x-L_0(3\mathsf x+4)\}$
and value $P^+(L_0,\mathsf x,\mathsf t,0)=L_0^3\{L_0(\mathsf x+4)-3\mathsf t\mathsf x\}$.
Since $\mathsf t\mathsf x=L_0\mathsf w<L_0$, the leading coefficient is
negative and the value at $\mathsf y=0$ exceeds $L_0^4(1+\mathsf x)>0$.
So $P^+$ is strictly concave in $\mathsf y$ on the admissible interval and
positive at its lower end; it is positive on the whole interval as soon as
it is positive at the upper end, which lies on $\mathsf S_V=0$ or on
$\mathsf S_Z=0$.

\emph{Step 4: the face $\mathsf S_Z=0$.}  There $ab=1$.  Put
$r=\mathsf y$, $h=\mathsf w$, $w_0=1-h\in(0,1)$, so that $b=gw_0$ and
$a=(gw_0)^{-1}$, and let $\mathcal T=g^2w_0^2\,(-Q)$; explicitly
\[
 \mathcal T=m\,(r-1)(rw_0^2-1)^2
 +r^2w_0^4-3r^2w_0^3+3r^2w_0^2-rw_0-r+1,
\]
so $\partial_m\mathcal T=(r-1)(rw_0^2-1)^2<0$ (note $rw_0^2=b^2<ab=1$).
The margin $\mathsf S_V\ge0$ says $m\le M=\{(1-r)-4rh^2\}/(1-r)^2$.  If
$M\ge1$, i.e. $r\le1-4h^2$ (so $h<1/2$), then
$\mathcal T(m)>\mathcal T(1)=rw_0(rw_0-1)(rw_0^2-2w_0+1)$, and the last
factor increases in $r$, so
$rw_0^2-2w_0+1\le(1-4h^2)(1-h)^2-2(1-h)+1=-h^2(2h-3)(2h-1)<0$; with
$rw_0<1$ this gives $\mathcal T(1)>0$.  If $M<1$, then
$\mathcal T(m)\ge\mathcal T(M)=rh\,\mathcal P(h,r)/(r-1)$ with
\[
 \mathcal P=-4h^5r^2+16h^4r^2-24h^3r^2+8h^3r+19h^2r^2-19h^2r
 -10hr^2+16hr-6h+3r^2-6r+3 ,
\]
whose leading coefficient in $r$ is $-(h-1)^2(4h^3-8h^2+4h-3)>0$, so
$\mathcal P$ is strictly convex in $r$; the conditions $0<M<1$ confine
$r$ to $\max(0,1-4h^2)<r<(1+4h^2)^{-1}$, and at the endpoints
$\mathcal P(h,(1+4h^2)^{-1})=-4h^4(17h+3)/(1+4h^2)^2<0$,
$\mathcal P(h,1-4h^2)=-4h^4(h-1)(2h-3)(2h-1)(4h^2-4h-1)<0$ for $h<1/2$,
and $\mathcal P(h,0)=3(1-2h)\le0$ for $h\ge1/2$.  Hence $\mathcal P<0$ on
the open interval, $\mathcal T(M)>0$, and $\mathcal T>0$, i.e. $-Q>0$,
throughout the face.

\emph{Step 5: the face $\mathsf S_V=0$.}  There $\mathsf v=\rho/2$.  Put
$r=\mathsf y$, $c=\mathsf w=\rho/(2g)$ and $z=ab$.  The face equation gives
$m=M_0/(1-r)^2$ with $M_0=1-r-4rc^2$, and $0<m<1$ makes $M_0>0$ and
$M_1=r+4c^2-1>0$; Step~1 reads
$z_0<z<1$ with $z_0=r(1-c)\{M_0+2rM_1c\}/M_0$.  With $(r,c)$ fixed the
cleared $-Q$ is a cubic in $z$.  Its first and third Bernstein control
points on $[z_0,1]$ have exact positive rational identities with $133$ and
$111$ terms; its second control point is regularized at the cusp
$r=1$, $c=0$ by the two charts $c=h/2$, $r=1-h^2+h^4u/(1+h^2)$ and
$c=(1+h)/2$, $r=u/(h^2+2h+2)$, which encode $M_0,M_1\ge0$ and leave the
margin $M_0(1-z_0)>0$, and has positive rational identities with $246$
and $165$ terms in the two charts; its fourth control point is the value
at $z=1$, which lies on $\mathsf S_Z=0$ and is positive by Step~4.  The
four identities, $655$ positive rational terms in all, are frozen in
\path{data/universal_right_sv_certificates.json}, which certifies the
strict positivity of the Bernstein control coefficients $B_0$, $B_1$ (in
both charts) and $B_2$ of the cleared $-Q$ cubic on $[z_0,1]$ under the
margins $M_0,M_1,M_0(1-z_0)>0$; the verifier
\path{scripts/universal_right_fold_proof.py} reconstructs them over
$\mathbb Q$ without an optimizer and also checks the identities of
Steps~1, 3 and~4.  Hence $-Q>0$ on the face, and by Step~3 on the whole
sheet, the $\alpha_1$ end included for $k<k_0$.
\end{proof}

\bibliographystyle{plain}
\bibliography{../references/references}

\end{document}
